%  LaTeX support: latex@mdpi.com 
%  For support, please attach all files needed for compiling as well as the log file, and specify your operating system, LaTeX version, and LaTeX editor.

%=================================================================
\documentclass[biology,article,submit,pdftex,oneauthor]{Definitions/mdpi}
%\documentclass[preprints,article,submit,pdftex,moreauthors]{Definitions/mdpi} 
% For posting an early version of this manuscript as a preprint, you may use "preprints" as the journal. Changing "submit" to "accept" before posting will remove line numbers.

%--------------------
% Class Options:
%--------------------
%----------
% journal
%----------
% Choose between the following MDPI journals:
% accountaudit, acoustics, actuators, addictions, adhesives, adolescents, aerobiology, aerospace, agriculture, agriengineering, agrochemicals, agronomy, ai, aichem, aieng, aimater, aimed, aipa, air, aisens, algorithms, allergies, alloys, amh, analog, analytica, analytics, anatomia, anesthres, animals, antibiotics, antibodies, antioxidants, applbiosci, appliedchem, appliedmath, appliedphys, applmech, applmicrobiol, applnano, applsci, aquacj, architecture, arm, arthropoda, asc, asi, astronautics, astronomy, atmosphere, atoms, audiolres, automation, axioms, bacteria, batteries, bdcc, beverages, biochem, bioengineering, biologics, biology, biomass, biomechanics, biomed, biomedicines, biomedinformatics, biomimetics, biomolecules, biophysica, bioresourbioprod, biosensors, biosphere, biotech, birds, blockchains, bloods, blsf, brainsci, breath, buildings, cancers, carbon, cardio, cardiogenetics, cardiovascmed, catalysts, cells, ceramics, challenges, chemengineering, chemistry, chemosensors, chemproc, children, chips, cimb, civileng, cleantechnol, climate, clinbioenerg, clinpract, clockssleep, cmd, cmtr, coasts, coatings, colloids, colorants, commodities, complexities, complications, compounds, computation, computers, condensedmatter, conservation, constrmater, cosmetics, covid, crops, cryo, cryptography, crystals, csmf, ctn, culture, curroncol, cyber, dairy, data, ddc, dentistry, dermato, dermatopathology, designs, devices, dhi, diabetology, diagnostics, dietetics, digital, disabilities, diseases, diversity, dna, drones, dynamics, earth, ebj, ecm, ecologies, edm, eesp, electricity, electrochem, electronicmat, electronics, encyclopedia, endocrines, energies, eng, engproc, entomology, entropy, environments, environremediat, epidemiologia, epigenomes, esa, est, fermentation, fibers, fintech, fire, fishes, fluids, foods, forecasting, forensicsci, forests, fossstud, foundations, fractalfract, fuels, future, futureinternet, futurepharmacol, futurephys, futuretransp, galaxies, gases, gastroent, gastrointestdisord, gastronomy, gels, genes, geographies, geohazards, geomatics, geometry, geosciences, geotechnics, geriatrics, germs, glacies, grasses, green, greenhealth, gucdd, hardware, hazardousmatters, healthcare, hearts, hemato, hematolrep, hep, heritage, higheredu, highthroughput, horticulturae, hospitals, hydrobiology, hydrogen, hydrology, hydropower, hygiene, idr, iic, ijem, ijerph, ijgi, ijmd, ijms, ijns, ijom, ijpb, ijt, ijtm, ijtpp, ime, immuno, informatics, information, infrastructures, inorganics, insects, instruments, inventions, iot, j, jaestheticmed, jal, jcdd, jcm, jcp, jcrm, jcs, jcto, jdad, jdb, jdream, jemr, jeta, jfb, jfmk, jgbg, jgg, jimaging, jlpea, jmahp, jmmp, jmms, jmp, jmse, jne, jnt, jof, joi, joitmc, joma, jop, joptm, jor, jox, jpbi, jphytomed, jpm, jsan, jtaer, jvd, jzbg, kidneydial, kinasesphosphatases, knowledge, labmed, laboratories, lae, land, life, lights, limnolrev, lipidology, liquids, livers, logics, logistics, lubricants, lymphatics, machines, macromol, magnetism, magnetochemistry, make, marinedrugs, materials, materproc, mathematics, mca, measurements, medicina, medicines, medsci, membranes, merits, metabolites, metals, meteorology, methane, metrics, metrology, micro, microarrays, microbiolres, microelectronics, micromachines, microorganisms, microplastics, microwave, minerals, mining, mmphys, modelling, molbank, molecules, mps, msf, mti, multimedia, muscles, nanoenergyadv, nanomanufacturing, nanomaterials, ncrna, ndt, network, neuroglia, neuroimaging, neurolint, neurosci, nitrogen, notspecified, nursrep, nutraceuticals, nutrients, obesities, occuphealth, oceans, ohbm, onco, optics, oral, organics, organoids, osteology, oxygen, pandemics, parasites, parasitologia, particles, pathogens, pathophysiology, pediatrrep, pets, pharmaceuticals, pharmaceutics, pharmacoepidemiology, pharmacy, philosophies, photochem, photonics, phycology, physchem, physics, physiologia, plants, plasma, platforms, pollutants, polymers, polysaccharides, populations, poultry, powders, precipitation, precisoncol, preprints, proceedings, processes, prosthesis, proteomes, psf, psychiatryint, psychoactives, purification, quantumrep, quaternary, qubs, radiation, rdt, reactions, realestate, receptors, recycling, regeneration, remotesensing, reports, reprodmed, resources, rheumato, rjpm, robotics, rsee, ruminants, safety, sci, scipharm, sclerosis, seeds, sensors, separations, sexes, shi, signals, sinusitis, siuj, skins, smartcities, sna, societies, software, soilsystems, solar, solids, spectroscj, sports, standards, stats, std, stratsediment, stresses, surfaces, surgeries, suschem, sustainability, symmetry, synbio, systems, tae, targets, taxonomy, technologies, telecom, test, textiles, thalassrep, therapeutics, thermo, timespace, tomography, toxics, toxins, tph, transplantology, transportation, traumacare, traumas, tri, tropicalmed, universe, urbansci, uro, vaccines, vehicles, venereology, vetsci, vibration, virtualworlds, viruses, vision, waste, water, welding, wem, wevj, wild, wind, women, world, zoonoticdis

%---------
% article
%---------
% The default type of manuscript is "article", but can be replaced by: 
% abstract, addendum, article, benchmark, book, bookreview, briefcommunication, briefreport, casereport, changes, clinicopathologicalchallenge, comment, commentary, communication, conceptpaper, conferenceproceedings, correction, conferencereport, creative, datadescriptor, discussion, entry, expressionofconcern, extendedabstract, editorial, essay, erratum, fieldguide, hypothesis, interestingimages, letter, meetingreport, monograph, newbookreceived, obituary, opinion, proceedingpaper, projectreport, reply, retraction, review, perspective, protocol, shortnote, studyprotocol, supfile, systematicreview, technicalnote, viewpoint, guidelines, registeredreport, tutorial,  giantsinurology, urologyaroundtheworld
% supfile = supplementary materials

%----------
% submit
%----------
% The class option "submit" will be changed to "accept" by the Editorial Office when the paper is accepted. This will only make changes to the frontpage (e.g., the logo of the journal will get visible), the headings, and the copyright information. Also, line numbering will be removed. Journal info and pagination for accepted papers will also be assigned by the Editorial Office.

%------------------
% moreauthors
%------------------
% If there is only one author the class option oneauthor should be used. Otherwise use the class option moreauthors.

%---------
% pdftex
%---------
% The option pdftex is for use with pdfLaTeX. Remove "pdftex" for (1) compiling with LaTeX & dvi2pdf (if eps figures are used) or for (2) compiling with XeLaTeX.

%=================================================================
% MDPI internal commands - do not modify
\firstpage{1} 
\makeatletter 
\setcounter{page}{\@firstpage} 
\makeatother
\pubvolume{1}
\issuenum{1}
\articlenumber{0}
\pubyear{2026}
\copyrightyear{2026}
%\externaleditor{Firstname Lastname} % More than 1 editor, please add `` and '' before the last editor name
\datereceived{ } 
\daterevised{ } % Comment out if no revised date
\dateaccepted{ } 
\datepublished{ } 
%\datecorrected{} % For corrected papers: "Corrected: XXX" date in the original paper.
%\dateretracted{} % For retracted papers: "Retracted: XXX" date in the original paper.
%\doinum{} % Used for some special journals, like molbank
%\pdfoutput=1 % Uncommented for upload to arXiv.org
%\CorrStatement{yes}  % For updates
%\longauthorlist{yes} % For many authors that exceed the left citation part
%\IsAssociation{yes} % For association journals

%=================================================================
% Add packages and commands here. The following packages are loaded in our class file: fontenc, inputenc, calc, indentfirst, fancyhdr, graphicx, epstopdf, lastpage, ifthen, float, amsmath, amssymb, lineno, setspace, enumitem, mathpazo, booktabs, titlesec, etoolbox, tabto, xcolor, colortbl, soul, multirow, microtype, tikz, totcount, changepage, attrib, upgreek, array, tabularx, pbox, ragged2e, tocloft, marginnote, marginfix, enotez, amsthm, natbib, hyperref, cleveref, scrextend, url, geometry, newfloat, caption, draftwatermark, seqsplit
% cleveref: load \crefname definitions after \begin{document}

\usepackage{longtable}
\usepackage{placeins}  % \FloatBarrier for section-level float control
\raggedbottom           % reduce stretched vertical whitespace on float-heavy pages
\setlength{\headheight}{20pt}

% --- Reduce float spacing for tighter layout ---
\setlength{\textfloatsep}{8pt plus 2pt minus 2pt}   % between float and text
\setlength{\floatsep}{8pt plus 2pt minus 2pt}       % between adjacent floats
\setlength{\intextsep}{8pt plus 2pt minus 2pt}      % for in-text floats

% --- Reduce caption spacing ---
\setlength{\abovecaptionskip}{4pt}
\setlength{\belowcaptionskip}{2pt}

% --- Reduce display-math spacing ---
\AtBeginDocument{%
  \setlength{\abovedisplayskip}{6pt plus 2pt minus 2pt}%
  \setlength{\belowdisplayskip}{6pt plus 2pt minus 2pt}%
  \setlength{\abovedisplayshortskip}{3pt plus 1pt minus 1pt}%
  \setlength{\belowdisplayshortskip}{3pt plus 1pt minus 1pt}%
}

% --- Tighter section spacing (override MDPI defaults: 0pt/12pt/3pt) ---
\titlespacing{\section}{0pt}{8pt}{2pt}
\titlespacing{\subsection}{0pt}{8pt}{2pt}
\titlespacing{\subsubsection}{0pt}{8pt}{2pt}
\titlespacing{\paragraph}{0pt}{6pt}{2pt}

%=================================================================
% Please use the following mathematics environments: Theorem, Lemma, Corollary, Proposition, Characterization, Property, Problem, Example, ExamplesandDefinitions, Hypothesis, Remark, Definition, Notation, Assumption
%% For proofs, please use the proof environment (the amsthm package is loaded by the MDPI class).

%=================================================================
% Full title of the paper (Capitalized)
\Title{CLOP-DiT: Structured-Metadata-Conditioned Single-Cell Latent Generation via Contrastive Language-Omics Pretraining and Diffusion Transformers}

% Author Orchid ID: enter ID or remove command
\newcommand{\orcidauthorA}{0009-0001-8329-0108} % Add \orcidA{} behind the author's name
%\newcommand{\orcidauthorB}{0000-0000-0000-000X} % Add \orcidB{} behind the author's name

% Authors, for the paper (add full first names)
\Author{Zeyu Fu $^{1,}$*\orcidA{}}

%\longauthorlist{yes}

% MDPI internal command: Authors, for metadata in PDF
\AuthorNames{Zeyu Fu}

% Affiliations / Addresses (from MDPI Biology 14/6/679 — scRL; Zeyu Fu affiliation)
\address[1]{%
$^{1}$ \quad State Key Laboratory of Trauma and Chemical Poisoning, Institute of Combined Injury, Chongqing Engineering Research Center for Nanomedicine, College of Preventive Medicine, Army Medical University, Chongqing 400038, China; zeyu.fu@tmmu.edu.cn}

% Contact information of the corresponding author
\corres{Correspondence: zeyu.fu@tmmu.edu.cn.}

% Current address and/or shared authorship
%\firstnote{Current address: Affiliation.}  
% Current address should not be the same as any items in the Affiliation section.

%\secondnote{These authors contributed equally to this work.}
% The commands \thirdnote{} till \eighthnote{} are available for further notes.

%\simplesumm{} % Simple summary

%\conference{} % An extended version of a conference paper

% Abstract (Do not insert blank lines, i.e. \\) 
\abstract{Generating synthetic single-cell states from structured biological descriptors could support simulation and hypothesis generation, but practical utility remains unproven. We present CLOP-DiT, a two-stage pipeline that combines Contrastive Language--Omics Pretraining (CLOP) with a conditional Diffusion Transformer (DiT) trained by flow matching. CLOP aligns BiomedBERT text embeddings and scGPT cell embeddings in a shared 512-dimensional space, and DiT then samples scGPT latents conditioned on a fixed five-field template comprising cell type, tissue, organism, marker genes, and disease context. The method should therefore be interpreted as structured-metadata conditioning rather than free-form natural-language generation. Across 69 deduplicated cell types from 80 GEO datasets (220,304 cells), CLOP expands inter-type separation from 0.006 to 0.778 in cosine distance, and CLOP-DiT achieves 36.9\% KNN accuracy and 81.0\% steering in a high-fidelity regime (CFG = 2.0, Euler solver). A high-diversity regime (CFG = 1.0, Midpoint solver) reaches a diversity ratio of 0.93 at 80.7\% steering, and three independent runs confirm stable steering and diversity estimates. A frozen scGPT decoder can map generated latents back to gene expression, but this mapping is many-to-one; accordingly, latent-space metrics are the primary indicators of generative quality. Decoder-reconstructed mean expression is highly concordant with real data (Pearson $r > 0.999$), yet this concordance primarily reflects the decoder's consistency rather than faithful generation: any latent near the training manifold decodes to a similar expression profile. Per-gene variance correlation is near zero and within-type gene--gene correlation structure is weakly preserved, indicating substantial loss of cell-to-cell heterogeneity. Downstream analyses show partial transfer into clustering, classifier transfer, and differential-expression workflows, but generated cells remain distinguishable from real data (discriminator AUC 0.656). External marker validation against PanglaoDB and CellMarker~2.0 supports the biological plausibility of the conditioning markers. A conditioning field ablation shows that removing marker genes from prompts reduces steering accuracy from 99.8\% to 62.4\% (metadata-only), and a swap-label permutation test confirms that generation follows the conditioning signal in 100\% of tested pairs, providing causal evidence against pure leakage. However, residual leakage cannot be fully excluded. A pilot rare-cell augmentation study was negative, and on the primary common-metrics-only benchmark the Gaussian baseline outperforms CLOP-DiT. All results are restricted to in-distribution human and mouse cancer/developmental datasets; no stringent out-of-distribution generalization test has been conducted. CLOP-DiT therefore establishes the feasibility of structured-metadata-conditioned single-cell latent generation, but not yet its practical utility for downstream biology.}

% Keywords
\keyword{single-cell RNA-seq; generative model; diffusion transformer; contrastive learning; structured-metadata conditioning; flow matching; latent generation; foundation model} 

% The fields PACS, MSC, and JEL may be left empty or commented out if not applicable
%\PACS{J0101}
%\MSC{}
%\JEL{}

%%%%%%%%%%%%%%%%%%%%%%%%%%%%%%%%%%%%%%%%%%
% Only for the journal Diversity
%\LSID{\url{http://}}

%%%%%%%%%%%%%%%%%%%%%%%%%%%%%%%%%%%%%%%%%%
% Only for the journal Applied Sciences
%\featuredapplication{Authors are encouraged to provide a concise description of the specific application or a potential application of the work. This section is not mandatory.}
%%%%%%%%%%%%%%%%%%%%%%%%%%%%%%%%%%%%%%%%%%

%%%%%%%%%%%%%%%%%%%%%%%%%%%%%%%%%%%%%%%%%%
% Only for the journal Data
%\dataset{DOI number or link to the deposited data set if the data set is published separately. If the data set shall be published as a supplement to this paper, this field will be filled by the journal editors. In this case, please submit the data set as a supplement.}
%\datasetlicense{License under which the data set is made available (CC0, CC-BY, CC-BY-SA, CC-BY-NC, etc.)}

%%%%%%%%%%%%%%%%%%%%%%%%%%%%%%%%%%%%%%%%%%
% Only for the journal BioTech, Fishes, Neuroimaging and Toxins
%\keycontribution{The breakthroughs or highlights of the manuscript. Authors can write one or two sentences to describe the most important part of the paper.}

%%%%%%%%%%%%%%%%%%%%%%%%%%%%%%%%%%%%%%%%%%
% Only for the journal Encyclopedia
%\encyclopediadef{For entry manuscripts only: please provide a brief overview of the entry title instead of an abstract.}

%%%%%%%%%%%%%%%%%%%%%%%%%%%%%%%%%%%%%%%%%%
% Different journals have different requirements. Please check the specific journal guidelines in the "Instructions for Authors" on the journal's official website.
%\addhighlights{yes}
%\renewcommand{\addhighlights}{%
%
%\noindent The goal is to increase the discoverability and readability of the article via search engines and other scholars. Highlights should not be a copy of the abstract, but a simple text allowing the reader to quickly and simplified find out what the article is about and what can be cited from it. Each of these parts should be devoted up to 2~bullet points.\vspace{3pt}\\
%\textbf{What are the main findings?}
% \begin{itemize}[labelsep=2.5mm,topsep=-3pt]
% \item First bullet.
% \item Second bullet.
% \end{itemize}\vspace{3pt}
%\textbf{What are the implications of the main findings?}
% \begin{itemize}[labelsep=2.5mm,topsep=-3pt]
% \item First bullet.
% \item Second bullet.
% \end{itemize}
%}

%%%%%%%%%%%%%%%%%%%%%%%%%%%%%%%%%%%%%%%%%%
\begin{document}

%%%%%%%%%%%%%%%%%%%%%%%%%%%%%%%%%%%%%%%%%%
\section{Introduction}

Single-cell RNA sequencing (scRNA-seq) has transformed the study of cellular heterogeneity by enabling transcriptome-scale profiling at single-cell resolution~\cite{ref-scrnaseq-review,ref-10x}. Yet generating synthetic single-cell states from structured biological descriptors remains difficult. A model that could condition on cell identity, tissue context, organism, marker genes, and disease state would be attractive for controlled simulation and hypothesis generation, but it should be judged by whether it preserves biologically meaningful structure rather than by the fluency of the textual prompt alone.

Most current single-cell generative models condition on categorical labels or perturbation metadata rather than richer multi-field descriptions. Variational autoencoders such as scVI~\cite{ref-scvi,ref-scvi-tools} and scGen~\cite{ref-scgen}, geometric generative models~\cite{ref-scdiff}, and foundation models such as Geneformer~\cite{ref-geneformer} and scBERT~\cite{ref-scbert,ref-scfoundation,ref-cellplm,ref-uce} broaden the modeling landscape, but they do not jointly learn a structured text--cell alignment space and a conditional diffusion-style generator. Cell2Sentence~\cite{ref-cell2sentence} and GenePT~\cite{ref-genept} move closer to language-based biology, yet they address different problem formulations and do not combine contrastive cross-modal alignment with latent flow matching. Table~\ref{tab:method-comparison} summarizes the key differences.

\begin{table}[!htbp]
\caption{Comparison of CLOP-DiT with related methods for conditional single-cell generation and annotation. ``Cond.\ modality'' indicates the conditioning input type; ``Gen.\ paradigm'' indicates the generation approach.\label{tab:method-comparison}}
\centering
\small
\begin{tabular}{lccc}
\toprule
\textbf{Method} & \textbf{Cond.\ Modality} & \textbf{Gen.\ Paradigm} & \textbf{Output} \\
\midrule
scVI~\cite{ref-scvi}           & Categorical (batch/type) & cVAE               & Expression \\
scGen~\cite{ref-scgen}         & Categorical (perturbation) & VAE + shift       & Expression \\
Geneformer~\cite{ref-geneformer} & Rank-value tokens & Masked LM           & Cell embeddings \\
Cell2Sentence~\cite{ref-cell2sentence} & Gene-rank text   & LLM decoding        & Gene-rank text \\
GenePT~\cite{ref-genept}       & Gene text embeddings     & None (annotation)   & Annotations \\
\textbf{CLOP-DiT (ours)}       & \textbf{Structured 5-field text} & \textbf{DiT + flow matching} & \textbf{Gene expression} \\
\bottomrule
\end{tabular}
\end{table}

Recent advances in contrastive multimodal learning and diffusion modeling motivate a direct bridge between biological descriptions and cell representations. CLIP/SigLIP-style objectives show that strong cross-modal alignment can emerge from contrastive training~\cite{ref-clip,ref-siglip}, while Diffusion Transformers and flow matching provide scalable conditional generation in continuous latent spaces~\cite{ref-dit,ref-flow-matching,ref-albergo-flow,ref-rectified-flow}. These ingredients suggest a natural design for structured biological conditioning: first build a well-separated text--cell condition space, then learn a conditional sampler in a biologically meaningful latent representation.

We present CLOP-DiT, a two-stage pipeline that aligns structured metadata with scGPT cell embeddings and then samples scGPT latents conditioned on that aligned representation. The input is a fixed five-field template---cell type, tissue, organism, marker genes, and disease context---so the method should be interpreted as structured-metadata conditioning rather than free-form language understanding. A frozen scGPT decoder can optionally map generated latents back to expression, but because this decoder is many-to-one, latent-space metrics are the primary indicators of generative quality.

The main contributions are fourfold. First, CLOP uses PrototypeSigLIP with ZCA-whitened BiomedBERT embeddings to align text and cell representations in a shared 512-dimensional space and substantially improve inter-type separability. Second, a 1D Diffusion Transformer performs conditional flow matching in that latent space to generate type-specific cell embeddings. Third, we evaluate the model on 69 deduplicated cell types from 80 GEO datasets using complementary latent, diversity, and downstream biological metrics. Fourth, we report both the strengths and the limits of the approach: structured conditioning produces clear above-chance type specificity, but variance preservation, covariance preservation, and rare-cell augmentation remain unresolved.

Across the main operating points, CLOP-DiT reaches 36.9\% KNN accuracy and 81.0\% steering in a high-fidelity regime, while a high-diversity regime achieves a diversity ratio of 0.93 at 80.7\% steering. Multi-seed replication shows stable steering and diversity, whereas decoder-level analyses reveal strong mean reconstruction but weak preservation of within-population heterogeneity. CLOP-DiT should therefore be viewed as a proof-of-concept structured-metadata-to-latent conditional sampler with a clear roadmap for improvement, rather than as a finished tool for downstream biological analysis.

%%%%%%%%%%%%%%%%%%%%%%%%%%%%%%%%%%%%%%%%%%
\section{Materials and Methods}

Figure~\ref{fig:architecture} provides an overview of the pipeline.
We organize the Materials and Methods in three parts:
(1) Dataset curation and preprocessing (\ref{sec:dataset});
(2) Model architecture and training---covering CLOP alignment (\ref{sec:clop}),
DiT generation (\ref{sec:dit}), and scGPT decoding (\ref{sec:decoding});
(3) Evaluation framework and metrics (\ref{sec:evaluation}).
Together, these components form a reproducible pipeline for text-conditioned single-cell generation.

\begin{figure}[!htbp]
\centering
\includegraphics[width=\textwidth]{figures/fig01_architecture.pdf}
\caption{Overview of the CLOP-DiT pipeline. \textbf{Stage~1 (CLOP)} aligns structured text descriptions and scGPT cell embeddings in a shared 512-dimensional space after BiomedBERT whitening and projection. \textbf{Stage~2 (DiT)} performs conditional flow matching in scGPT latent space and generates a target latent from Gaussian noise using classifier-free guidance at inference. \textbf{Stage~3 (Decoding)} maps the generated latent back to gene expression through the frozen scGPT decoder.\label{fig:architecture}}
\end{figure}

\subsection{Dataset}\label{sec:dataset}

We curated 80 datasets from the Gene Expression Omnibus (GEO) spanning cancer and developmental biology, comprising 220,304 cells after quality control and highly variable gene selection~\cite{ref-seurat,ref-hvg} (2,000 highly variable genes per dataset, selected by mean expression and dispersion). Each cell was encoded into a 512-dimensional embedding by a pre-trained scGPT model. Cell-type annotations were organized into 1,088 unique text groups, each described by a structured text caption (not free-form natural language) incorporating cell type, marker genes, tissue of origin, organism, and disease context. Text descriptions were generated for each cell type using a structured template:
\begin{quote}\small\ttfamily
\{cell type\}, tissue: \{tissue\}, organism: \{organism\},\\
markers: \{top 5 DE marker genes\}, context: \{disease/condition\}
\end{quote}
\noindent where marker genes were identified per cell type by one-vs.-rest Wilcoxon rank-sum test~\cite{ref-wilcoxon} on the training set. As an independent check on marker gene circularity, we validated these training-set-derived markers against PanglaoDB~\cite{ref-panglao} and CellMarker~2.0~\cite{ref-cellmarker}; the high Jaccard overlap (0.39--0.45) indicates that the selected markers reflect biological consensus (see Discussion). The exact template and marker gene sources are available in {\small\texttt{src/data\_pipeline/\allowbreak{}subcluster\_annotation.py}} (caption generation) and {\small\texttt{scripts/data\_prep/\allowbreak{}02b\_enrich\_descriptions.py}} (evidence enrichment).

The corpus was split into 72 training datasets (199,670 cells) and 8 held-out validation datasets (20,634 cells). The 8 validation datasets were selected by study (not by cell) to ensure no sample overlap between training and validation splits, and were chosen to span the major tissue and cancer categories in the corpus (lung, gastric, skin, liver, blood) via stratified selection across tissue types; the full list of GEO accession identifiers for all 80 datasets is provided in Supplementary Table~S1, and the validation study identifiers are in Supplementary Table~S2. Within the 72 training datasets, a cell-type-stratified split reserves $\sim$10\% of cells per type for training-time validation, ensuring that all 69 deduplicated cell types are represented in both partitions. This held-out design tests interpolation within the training distribution; the 8 held-out datasets share tissue types and cell populations with the training corpus and do not represent a stringent out-of-distribution generalization test.

The relationship between the 1,088 text groups and the 69 evaluated cell types is as follows. The 1,088 text groups correspond to sub-cluster-level descriptions: each Leiden~\cite{ref-leiden} cluster within each dataset receives a unique text caption. Many sub-clusters share the same canonical cell type (e.g., ``CD8+ T cells'' in lung and liver datasets produce separate text groups). For evaluation, text groups are consolidated by core cell-type identity via a deduplication pipeline that (a)~removes uncharacterized or ambiguous sub-clusters (mapped to null in the annotation pipeline) and (b)~merges cross-dataset sub-clusters into canonical types, yielding 69 evaluable cell types after deduplication. This consolidation is implemented in \texttt{scripts/analysis/deduplicate\_captions.py} and produces the deduplicated cache (\texttt{data/cached\_latents\_v5.2/METADATA\_DEDUP.json}: 1,088 original captions $\to$ 69 deduplicated types, 167,245 cells). KNN accuracy is computed over these 69 consolidated types (random chance = $1/69 \approx 1.45\%$). All headline generation metrics, bootstrap confidence intervals, and benchmark comparisons operate on this consistent 69-type evaluation set.

The 80 datasets comprise both \textit{Homo sapiens} (59 datasets) and \textit{Mus musculus} (21 datasets)~\cite{ref-tabula-muris}, predominantly from tumor microenvironment and developmental contexts. Consequently, all generation and evaluation results should be interpreted as applying to human and mouse cancer and developmental single-cell biology; generalization to other organisms, non-cancer tissues, or perturbation conditions has not been tested.

Training and evaluation were implemented in Python~3.10 using PyTorch~2.1 (CUDA~12.0), scGPT~v0.2.1, and Hugging Face Transformers~4.36 for BiomedBERT; the complete environment specification is provided in \texttt{requirements.txt} and \texttt{configs/clop.yaml}.

\paragraph{Hardware and compute budget.}
All training and evaluation were conducted on a single NVIDIA GeForce RTX~5090 Laptop GPU (24\,GB VRAM) using PyTorch~2.1 with CUDA~12.0 and automatic mixed precision (AMP/FP16). The modest compute requirements ($<$2~GPU-hours total, $<$24~GB VRAM peak) mean that any Ampere-or-later GPU with $\geq$16~GB VRAM should suffice for reproduction. CLOP alignment training (60 epochs, batch size 1024) completes in approximately 10~minutes ($\sim$10\,s/epoch). DiT flow-matching training (200 epochs, batch size 1024) completes in approximately 75~minutes ($\sim$23\,s/epoch). End-to-end figure reproduction from cached embeddings (\texttt{scripts/pipeline/regenerate\_report.sh}) takes under 30~minutes. Total compute for a single full training run (CLOP + DiT + evaluation + visualization) is under 2~GPU-hours; no multi-GPU or distributed training was used. Inference for 69 cell types ($\sim$6,900 generated cells) takes approximately 2~minutes with 10-step Euler integration.

\subsection{CLOP: Contrastive Language--Omics Pretraining}\label{sec:clop}

CLOP aligns structured metadata descriptions and cell profiles into a shared 512-dimensional embedding space using a PrototypeSigLIP contrastive loss, producing the well-separated condition space required by the DiT. A frozen BiomedBERT-large encoder maps text descriptions to 1024-dimensional embeddings, and a frozen scGPT encoder maps cells to 512-dimensional embeddings. Dual three-layer MLP projectors (text: $1024 \to 1024 \to 1024 \to 512$; cell: $512 \to 512 \to 512 \to 512$) with BatchNorm~\cite{ref-batchnorm}, GELU activation~\cite{ref-gelu}, and dropout~\cite{ref-dropout} (0.2) project both modalities onto the $L_2$-normalized unit hypersphere.

\paragraph{ZCA whitening.}
Prior to projection, Zero-phase Component Analysis (ZCA) whitening is applied to the BiomedBERT embeddings to decorrelate dimensions and mitigate the collapse typical of frozen BERT representations. The transform is fitted once on the training text embeddings using the full $1024 \times 1024$ covariance matrix with regularisation $\epsilon = 10^{-4}$, then reused unchanged for validation and inference text. All components are retained; the goal is decorrelation rather than dimensionality reduction. In preliminary sensitivity checks, ZCA produced substantially lower pairwise cosine similarity than mean-centering or LayerNorm, and it was therefore adopted as the production preprocessing step.

The CLOP loss is:

\begin{linenomath}
\begin{equation}
\mathcal{L}_\text{CLOP} = \frac{1}{2}\left[\text{CE}\!\left(\frac{\mathbf{T}\mathbf{C}^\top}{\tau},\,\text{arange}(B)\right) + \text{CE}\!\left(\frac{\mathbf{C}\mathbf{T}^\top}{\tau},\,\text{arange}(B)\right)\right]
\label{eq:clop}
\end{equation}
\end{linenomath}

where $\mathbf{T}, \mathbf{C} \in \mathbb{R}^{B \times 512}$ are $L_2$-normalized projection matrices (each row is a projected text or cell embedding, respectively), $B$ is the batch size, $\mathbf{T}\mathbf{C}^\top \in \mathbb{R}^{B \times B}$ is the pairwise cosine similarity matrix, $\tau$ is a logit-scale parameter (following the CLIP convention where $\tau$ multiplies rather than divides the similarity matrix), and label smoothing $\epsilon = 0.1$ is applied. The targets $\text{arange}(B) = [0, 1, \ldots, B{-}1]$ enforce identity matching: each text embedding should match its corresponding cell embedding in the batch. In the production configuration, $\tau = 14.0$ is \textbf{fixed} (effective softmax temperature $1/\tau \approx 0.071$); the value was selected via grid search over $\{7.0, 10.0, 14.0, 20.0\}$ to maximise inter-type cosine separation on the validation split, yielding the sharpest condition space before training-accuracy collapse. The ablation study (Table~\ref{tab:ablation}) was conducted on a prior baseline (v8.2) that used a \emph{learnable} temperature; the ``Fixed temperature'' ablation row therein shows that switching from learned to fixed temperature reduced prototype accuracy by $-4.61$~pp under that earlier configuration. The production model (v9.3) subsequently adopted fixed temperature alongside a stratified validation split and updated regularization schedule, under which the fixed-temperature setting achieved satisfactory downstream generation quality (Section~\ref{sec:results}). The loss function follows the PrototypeSigLIP formulation, which aligns text prototypes to cell-group centroids rather than individual cells, with a cohesion regularization weight of 0.15. The CLOP aligner has 3.02M trainable parameters and is trained for 60 epochs with AdamW~\cite{ref-adamw} (lr~$= 5 \times 10^{-4}$, weight decay 0.06, cosine schedule~\cite{ref-cosine-schedule} with 5-epoch warmup).

The critical output of CLOP is the projected condition space. Raw scGPT cell-type centroids have a pairwise cosine similarity of 0.994 (cell types differ by only 0.6\%), whereas CLOP-projected conditions have a pairwise cosine of 0.222, amplifying inter-group separation by a factor of $\sim$130$\times$. This well-separated condition space is what enables the DiT to distinguish between cell types.

\subsection{DiT: Conditional Flow Matching with Diffusion Transformer}\label{sec:dit}

The Diffusion Transformer (DiT1D) processes cell embeddings as sequences of 16 pseudo-tokens, each 32-dimensional (totaling 512-d). The architecture comprises 8 AdaLN-Zero transformer blocks with 8-head self-attention (head dimension 64) and feed-forward layers ($512 \to 2048 \to 512$). Timestep information is encoded via sinusoidal embeddings projected to 512 dimensions, and CLOP condition vectors are projected from 512 to 512 dimensions. The combined condition $c_\text{combined} = t_\text{emb} + c_\text{emb}$ (element-wise summation, not concatenation) modulates each block through 6 AdaLN-Zero parameters ($\gamma_1, \beta_1, \alpha_1, \gamma_2, \beta_2, \alpha_2$), all zero-initialized for stable early training. The model has 22.10M parameters.

Training uses a flow matching objective~\cite{ref-flow-matching,ref-ot-cfm}:

\begin{linenomath}
\begin{equation}
\mathcal{L}_\text{FM} = \| v_\theta(z_t, t, c) - (z_1 - z_0) \|^2
\label{eq:fm}
\end{equation}
\end{linenomath}

where $z_0 \sim \mathcal{N}(0, I)$, $z_1$ is the real cell embedding, $z_t = (1-t)z_0 + tz_1$, and $t$ is drawn from a logit-normal distribution ($\mu{=}0, \sigma{=}1$), which concentrates training on intermediate timesteps where the velocity field is most complex~\cite{ref-sd3}. During training, classifier-free guidance (CFG)~\cite{ref-cfg} is enabled by dropping the condition with 15\% probability; when the condition is dropped, $c_\text{emb}$ is replaced by a zero vector of the same dimension ($\mathbf{0} \in \mathbb{R}^{512}$), serving as the unconditional null token. At inference:

\begin{linenomath}
\begin{equation}
v_\text{guided} = v_\text{uncond} + s \cdot (v_\text{cond} - v_\text{uncond})
\label{eq:cfg}
\end{equation}
\end{linenomath}

where $s$ is the guidance scale. DiT is trained for 200 epochs with EMA~\cite{ref-ema} (decay = 0.9999).

\subsection{Decoding via scGPT}\label{sec:decoding}

Generated latent vectors $z_1 \in \mathbb{R}^{512}$ are mapped back to gene expression using the frozen scGPT decoder. This step is useful for downstream biological inspection, but it introduces an important limitation: the decoder is many-to-one, so diverse latent vectors near the training manifold can decode to similar expression profiles. We therefore treat latent-space metrics---KNN accuracy, steering, diversity ratio, and centroid cosine---as the primary indicators of generative quality, whereas expression-space analyses should be interpreted as decoder-induced compatibility checks rather than direct measures of generation fidelity. Breaking this bottleneck will require a trainable decoder or lightweight adapter, which is left to future work.

\subsection{Evaluation Framework}\label{sec:evaluation}

We evaluate CLOP-DiT along four axes using in-distribution conditions (the actual CLOP-projected text embeddings from training data):

\begin{enumerate}
\item \textbf{KNN classification accuracy}: For each generated cell, a $k$-nearest-neighbor classifier ($k{=}15$, cosine metric, PCA-50 features) trained on 80\% of real data predicts its cell type. PCA is fitted on the real training split only and applied identically to generated embeddings; the same PCA transform is used across all method comparisons. Random chance is $1/69 \approx 1.45\%$. Because real cell-type centroids in scGPT space have pairwise cosine similarity of 0.994, small embedding perturbations can substantially affect KNN accuracy; the biological meaning of absolute KNN values should therefore be interpreted cautiously.
\item \textbf{Steering accuracy}: For random cell-type pairs $(A, B)$, we generate cells conditioned on $A$ and test whether the generated centroid is closer to the real cluster $A$ than $B$ in centered space. Random chance is 50\%. \textbf{Note:} At intermediate-to-high CFG values (e.g., CFG$\,\geq\,$1.5) with multi-step solvers, steering accuracy can saturate at 1.0 with bootstrap CI [1.0, 1.0], indicating a ceiling effect; this metric loses discriminative power once guidance is sufficient.
\item \textbf{Diversity ratio}: Ratio of generated within-group variance to real within-group variance. Ideal is 1.0; values $<$1 indicate mode sharpening; values $>$1 indicate excess noise.
\item \textbf{Coverage and density}: Coverage measures the fraction of real samples that have at least one generated neighbor within the real data's $k$-NN radius ($k{=}5$); it quantifies mode coverage. Density measures the average number of generated samples falling within each real sample's $k$-NN radius, normalised by $k$; it quantifies precision. Both are computed in PCA-50 embedding space.
\item \textbf{Centroid cosine similarity}: The cosine similarity between the real and generated cell-type centroids, computed in the 512-dimensional scGPT embedding space (uncentered; $\text{cos}(\boldsymbol{\mu}_\text{real}, \boldsymbol{\mu}_\text{gen})$). Reported per-type and averaged across 69 types.
\item \textbf{Downstream biological concordance}: Clustering alignment (ARI~\cite{ref-ari}, NMI~\cite{ref-nmi}), classifier transfer, and differential expression concordance between real and generated data.
\end{enumerate}

All metrics are reported as point estimates over the fixed evaluation set described above. Two configurations are designated as \emph{primary operating points}: CFG\,=\,2.0 with 10-step Euler integration (high-fidelity regime) and CFG\,=\,1.0 with 10-step Midpoint integration (high-diversity regime). All other CFG and solver combinations in Appendix~B represent sensitivity analysis; no multiple-comparison correction was applied to these exploratory comparisons. Bootstrap 95\% confidence intervals for the composite benchmark scores are shown in Figure~\ref{fig:benchmark}d; variability across single-measurement metrics (KNN, steering, diversity ratio) is characterized by the CFG sweep (Appendix~B, Table~\ref{tab:cfg-sweep}).

\textbf{Caveat on Fr\'{e}chet Distance (FD).} FD~\cite{ref-fid} compares distributional moments (mean and covariance) between real and generated samples. While widely used for unconditional evaluation, FD can be misleading for \emph{conditional} generation: a model that reproduces the global marginal distribution but ignores conditioning achieves a low FD without having learned the target conditional. FD is therefore reported as a secondary metric throughout; the primary metrics (KNN accuracy, steering, diversity ratio) are specifically designed for conditional evaluation.

The biological evidence is organized by figure family: Figures~\ref{fig:markers}--\ref{fig:expr-analysis} establish marker-level and gene-level fidelity; Figures~\ref{fig:diversity}--\ref{fig:noise-tradeoff} and~\ref{fig:expr-diversity} provide robustness readouts for guidance and diversity; Figures~\ref{fig:baselines}--\ref{fig:benchmark} report benchmarking; Figures~\ref{fig:downstream-pq} and~\ref{fig:downstream-r} test downstream biology. Figure~\ref{fig:eval-pipeline} provides a schematic overview of the complete evaluation pipeline, including data splitting, PCA fitting, KNN training, and metric aggregation stages.

\begin{figure}[!htbp]
\centering
\includegraphics[width=0.95\textwidth]{figures/fig02_evaluation_pipeline.pdf}
\caption{Evaluation pipeline schematic. Data flow from raw datasets through deduplication, stratified splitting, scGPT encoding, PCA fitting (on real training split only), KNN classifier training, and metric computation. The common-metrics-only composite (9 shared distributional metrics, marked in red) is the primary benchmark; the full 17-metric composite includes 8 downstream biology metrics reported only by CLOP-DiT and is structurally biased. Bootstrap confidence intervals ($B{=}1{,}000$) are computed by resampling the 69 evaluation types.\label{fig:eval-pipeline}}
\end{figure}

Having established the methodology and evaluation framework, we now present the empirical results.

\subsection{Statistical Analysis}\label{sec:stats}

\textbf{Primary endpoints.} Two operating points were pre-specified as primary comparisons: CFG\,=\,2.0 with 10-step Euler integration (high-fidelity regime, selected because KNN accuracy plateaus at CFG 2.0--3.0 in Table~\ref{tab:cfg-sweep}; the lower end was chosen to minimize diversity loss) and CFG\,=\,1.0 with 10-step Midpoint integration (high-diversity regime, selected to match DivR\,$\approx$\,1.0). The primary metrics are KNN top-1 classification accuracy (69-class, $k{=}15$, cosine metric on PCA-50 features) and steering accuracy (pairwise directional test, 1000 random pairs).

\textbf{Exploratory analyses.} All other CFG/solver/step combinations reported in Table~\ref{tab:cfg-sweep} (Appendix~B) constitute exploratory sensitivity analyses. No correction for multiple comparisons (e.g., Bonferroni, Holm) was applied to these comparisons. Rankings and performance claims refer only to the two primary operating points unless explicitly stated otherwise.

\textbf{Uncertainty quantification.} Bootstrap 95\% confidence intervals (percentile method, $B{=}1000$ resamples, seed 42) are reported for all headline metrics where per-unit (per-cell or per-type) data are available (see Results). These CIs capture evaluation sampling variability but not training-run variance, as all results derive from a single training run (seed 42). Formal superiority tests (e.g., permutation tests) were not conducted; all comparisons are descriptive.

\textbf{Multi-seed replication.} To characterise inter-run variance, we trained three independent CLOP + DiT pipeline replicas (seeds 42, 123, 456) with identical hyperparameters and report mean $\pm$ standard deviation for all core metrics in Table~\ref{tab:multi-seed} (see Section~\ref{sec:multi-seed-results}).

\textbf{Limitations of the statistical design.} (1)~No formal pre-registration: operating points were selected based on interpretable criteria post-hoc rather than by pre-registered protocol; this study is therefore exploratory rather than confirmatory. (2)~No multiple-comparison correction for the exploratory CFG sweep.

\subsection{Composite Score Formulation}\label{sec:composite}

The composite benchmark score aggregates all active metrics into a single ranking index. For each metric~$m$ with direction $d_m \in \{\text{lower}, \text{higher}\}$, let $v_{m,j}$ denote the raw value for method~$j$. We compute a min--max normalised score:
\begin{linenomath}
\begin{equation}
  s_{m,j} =
  \begin{cases}
    1 - \dfrac{v_{m,j} - \min_k v_{m,k}}{\max_k v_{m,k} - \min_k v_{m,k}} & \text{if } d_m = \text{lower},\\[6pt]
    \dfrac{v_{m,j} - \min_k v_{m,k}}{\max_k v_{m,k} - \min_k v_{m,k}} & \text{if } d_m = \text{higher},
  \end{cases}
  \label{eq:norm}
\end{equation}
\end{linenomath}
where $k$ ranges over all methods that report a non-null value for~$m$; methods lacking metric~$m$ receive $s_{m,j} = 0$. The composite score is the unweighted arithmetic mean over all~$M$ active metrics:
\begin{linenomath}
\begin{equation}
  C_j = \frac{1}{M}\sum_{m=1}^{M} s_{m,j}.
  \label{eq:composite}
\end{equation}
\end{linenomath}
Because not all methods support downstream biology (Section~\ref{sec:evaluation}, axis~6), they receive $s = 0$ on those metrics, which inflates the score of methods that do. In the current benchmark, 8 of 17 active metrics are reported only by CLOP-DiT; the remaining 9 distributional metrics are shared by all methods. \textbf{To enable fair like-for-like comparison, the common-metrics-only composite ($C_j^{\text{common}}$, restricted to the 9 shared distributional metrics) serves as the primary benchmark figure.} The full composite ($C_j$ over all $M = 17$ metrics) is reported as a supplementary capability ranking but should not be used for method comparison, as the 8 exclusive metrics structurally bias it in favour of CLOP-DiT. Table~\ref{tab:core-results} and the per-metric panels in Figure~\ref{fig:benchmark} report individual metrics for transparent comparison. Bootstrap 95\% confidence intervals (percentile method, $B{=}1{,}000$) for all headline generation metrics are provided in the supplementary materials (Table~S5).

%%%%%%%%%%%%%%%%%%%%%%%%%%%%%%%%%%%%%%%%%%
\section{Results}\label{sec:results}

We report in-distribution results along the four evaluation axes defined in the Methods, organized by evidence type:
training dynamics and embedding space (Figures~\ref{fig:training}--\ref{fig:embedding-space});
core generation quality metrics (Figures~\ref{fig:metrics}--\ref{fig:per-type-gen} and~\ref{fig:text-cell-alignment});
gene-level fidelity (Figures~\ref{fig:markers}--\ref{fig:expr-analysis});
conditioning robustness and diversity trade-offs (Figures~\ref{fig:cond-umap}--\ref{fig:noise-tradeoff} and~\ref{fig:expr-diversity});
baseline comparisons and benchmarking (Figures~\ref{fig:baselines}--\ref{fig:benchmark});
and downstream biological validation (Figures~\ref{fig:downstream-pq}--\ref{fig:downstream-r}).
This structure allows readers to assess evidence quality at multiple resolutions, from training convergence through partial downstream biological utility.

\subsection{Training Dynamics}

Figure~\ref{fig:training} shows the training curves for both stages of the pipeline. CLOP (top row) converges to a training loss of 1.187 at epoch 60, with batch training accuracy reaching 87\%. The validation accuracy of $\sim$2.5\% appears modest but represents 6.4$\times$ random chance (1/256 = 0.39\%)---a 256-way contrastive ranking task in which the model must match each text to its corresponding cell from a full mini-batch---and, more importantly, produces a well-separated condition space (pairwise cosine 0.222 vs.\ 0.994 for raw embeddings). DiT (bottom row) achieves a final validation velocity cosine of 0.990 with EMA, indicating near-perfect velocity field prediction. Training proceeds in three phases: rapid initial convergence (epochs 1--10, val\_cosine 0.13 $\to$ 0.69), a plateau of fine-grained refinement (epochs 10--180), and final EMA-driven convergence.

\begin{figure}[!htbp]
\centering
\includegraphics[width=0.95\textwidth]{figures/fig03_training_dynamics.pdf}
\caption{Training dynamics of the CLOP-DiT pipeline. (\textbf{a--d})~CLOP loss, fixed logit scale, batch accuracy, and embedding-quality metrics over 60 epochs. (\textbf{e--h})~DiT train/validation loss, velocity cosine similarity, learning-rate schedule, and normalized convergence comparison across training progress. Together, the panels show stable CLOP separation and near-converged DiT velocity prediction.\label{fig:training}}
\end{figure}

\subsection{CLOP Embedding Space}

Figure~\ref{fig:embedding-space} visualizes the embedding space at two levels. The top row shows the CLOP-projected embedding space via UMAP~\cite{ref-umap}: text and cell embeddings jointly projected into the shared 512-dimensional space form distinct clusters, with text embeddings co-localizing with their corresponding cell-type clusters, confirming successful cross-modal alignment. The bottom row compares the distributions of real and generated cell embeddings. Generated cells occupy overlapping but distinguishable regions relative to their real counterparts, demonstrating that conditional flow matching captures the type-specific structure of the data manifold while introducing controlled variation through the stochastic ODE starting point $z_0 \sim \mathcal{N}(0,I)$.

\begin{figure}[!htbp]
\centering
\includegraphics[width=0.95\textwidth]{figures/fig04_embedding_space.pdf}
\caption{Embedding space analysis. (\textbf{a--c})~CLOP-projected text and cell embeddings form aligned clusters in the shared space. (\textbf{d--f})~Real and generated scGPT embeddings occupy overlapping but distinguishable regions, indicating that conditional generation follows the coarse structure of the data manifold without fully matching real-cell geometry.\label{fig:embedding-space}}
\end{figure}

\subsection{Core Evaluation Metrics}

Table~\ref{tab:core-results} and Figure~\ref{fig:metrics} summarize the core results across generation configurations. At CFG = 2.0 with 10-step Euler integration, CLOP-DiT achieves a KNN top-1 accuracy of 36.9\% ($25\times$ random chance of 1.45\%; 52 points below the 89\% real-data baseline), KNN top-5 of 55.3\%, steering accuracy of 81.0\%, and linear classifier accuracy of 51.1\%. Note that the Fr\'{e}chet distance (FD) reported in Table~\ref{tab:core-results} is often misleading for conditional generation, because global mean/covariance matching tends to favour unconditional baselines; FD should therefore be interpreted alongside the directional metrics (KNN, steering, diversity). Unconditional generation ($s{=}0$) collapses to random chance (KNN = 1.0\%, steering = 47.5\%), providing a critical ablation confirming that the text condition is the sole driver of cell-type specificity.

\begin{table}[!htbp]
\caption{Core evaluation results. KNN-1 and KNN-5: top-1 and top-5 $k$-nearest-neighbor accuracy among 69 deduplicated cell types (random = 1.45\%). Steering: pairwise directional accuracy (random = 50\%). DivR: diversity ratio with ideal = 1.0. LinAcc: logistic regression accuracy in embedding space (trained on real, evaluated on generated; distinct from the expression-space downstream classifier in Figure~\ref{fig:downstream-pq}). FD: Fr\'{e}chet distance~\cite{ref-fid} (lower is better, but misleading for conditional generation; see text).\label{tab:core-results}}
\begin{tabularx}{\textwidth}{lCCCCCC}
\toprule
\textbf{Method} & \textbf{KNN-1} & \textbf{KNN-5} & \textbf{Steering} & \textbf{DivR} & \textbf{LinAcc} & \textbf{FD} \\
\midrule
Real data        & 0.890 & 0.993 & ---   & 1.000 & 0.942 & 0.00 \\
CFG=2.0 Euler-10 & 0.369 & 0.553 & 0.810 & 0.513 & 0.511 & 3.17 \\
CFG=3.0 Euler-10 & 0.367 & 0.554 & 0.802 & 0.473 & 0.508 & 3.46 \\
CFG=1.0 Mid-10   & 0.288 & 0.461 & 0.807 & 0.929 & 0.357 & 2.64 \\
CFG=1.0 Euler-50 & 0.293 & 0.466 & 0.807 & 0.919 & 0.365 & 2.65 \\
CFG=0.0 (uncond) & 0.010 & 0.051 & 0.475 & 1.833 & 0.010 & 0.58 \\
Gaussian          & 0.011 & 0.048 & 0.466 & 2.277 & 0.009 & 0.03 \\
\midrule
Random chance    & 0.010 & ---   & 0.500 & ---   & ---   & ---  \\
\bottomrule
\end{tabularx}
\end{table}

\begin{figure}[!htbp]
\centering
\includegraphics[width=0.95\textwidth]{figures/fig05_metrics_summary.pdf}
\caption{Metrics dashboard summarizing conditional generation quality. (\textbf{a})~Training convergence summary. (\textbf{b})~Multi-metric comparison between CLOP-DiT and the Gaussian baseline, including inverse FD, coverage, centroid cosine, diversity ratio, and per-gene Pearson correlation. (\textbf{c})~Diversity health gauges. (\textbf{d})~Expression-fidelity summary across configurations.\label{fig:metrics}}
\end{figure}

\textbf{Bootstrap confidence intervals.} To quantify evaluation sampling uncertainty, we computed bootstrap 95\% CIs by resampling the 69 cell types in the generation cache with replacement ($B{=}1{,}000$, percentile method). All 69 deduplicated types have both real and generated cells, so the bootstrap covers the full evaluation set consistently. At an intermediate reference configuration (CFG = 1.5, Midpoint, 20 steps; selected to balance fidelity and diversity between the two primary operating points): KNN top-1 = 0.509 [0.421, 0.586], KNN top-5 = 0.728 [0.658, 0.794], steering = 1.000 [1.000, 1.000] (\textbf{ceiling effect}: all bootstrap resamples achieve perfect pairwise accuracy at this configuration, indicating that the steering metric saturates at moderate-to-high CFG and becomes non-discriminative; this ceiling should be considered when interpreting steering differences across configurations), diversity ratio = 0.969 [0.963, 0.975], linear accuracy = 0.583 [0.518, 0.646], and centroid cosine = 0.898 [0.875, 0.913]. The narrow diversity ratio and centroid cosine CIs reflect uniformly high performance across types; the wider KNN CIs reflect type-dependent classification difficulty. These CIs capture evaluation sampling variability; inter-run variance across three training seeds is reported separately in Table~\ref{tab:multi-seed}.

\subsection{Multi-Seed Validation}\label{sec:multi-seed-results}

To quantify inter-run variance from stochastic initialization, data shuffling, and mini-batch ordering, we trained three independent CLOP + DiT pipelines (seeds 42, 123, 456) with identical hyperparameters and evaluated each under both operating regimes (Table~\ref{tab:multi-seed}).

\begin{table}[!htbp]
\caption{Multi-seed validation (3 independent training runs). Mean $\pm$ standard deviation for core metrics under both operating regimes. All three seeds use identical hyperparameters (Table~\ref{tab:hyperparams}), identical data splits, and identical evaluation protocol.\label{tab:multi-seed}}
\centering
\small
\begin{tabular}{lcccc}
\toprule
 & \multicolumn{2}{c}{\textbf{High-Fidelity (CFG=2.0, Euler)}} & \multicolumn{2}{c}{\textbf{High-Diversity (CFG=1.0, Midpoint)}} \\
\cmidrule(lr){2-3}\cmidrule(lr){4-5}
\textbf{Metric} & \textbf{Mean $\pm$ Std} & \textbf{Range} & \textbf{Mean $\pm$ Std} & \textbf{Range} \\
\midrule
KNN Top-1      & $0.348 \pm 0.082$ & [0.258, 0.419] & $0.273 \pm 0.048$ & [0.220, 0.312] \\
KNN Top-5      & $0.535 \pm 0.016$ & [0.524, 0.553] & $0.465 \pm 0.038$ & [0.430, 0.504] \\
Steering       & $0.808 \pm 0.009$ & [0.798, 0.817] & $0.801 \pm 0.008$ & [0.792, 0.804] \\
Diversity Ratio & $0.511 \pm 0.015$ & [0.496, 0.525] & $0.924 \pm 0.016$ & [0.906, 0.937] \\
Linear Acc.    & $0.488 \pm 0.023$ & [0.465, 0.511] & $0.355 \pm 0.097$ & [0.257, 0.452] \\
Centroid Cos.  & $0.892 \pm 0.012$ & [0.878, 0.900] & $0.892 \pm 0.010$ & [0.880, 0.898] \\
\bottomrule
\end{tabular}
\end{table}

Inter-run standard deviations are small relative to effect sizes: steering accuracy varies by $< 1$ percentage point ($\sigma = 0.009$), diversity ratio by $< 2$ percentage points ($\sigma = 0.015$--$0.016$), and centroid cosine by $\sim 1\%$ ($\sigma = 0.010$--$0.012$). KNN top-1 accuracy exhibits the largest inter-seed variance ($\sigma = 0.082$ in the high-fidelity regime), reflecting the sensitivity of nearest-neighbor classification to the decision boundary in the tightly-packed scGPT embedding space. Critically, the qualitative conclusions---$25\times$ above random chance, $\sim$52-point gap from real data, near-ideal diversity at CFG\,=\,1.0---are stable across all three seeds.

\subsection{Per-Type Generation Quality and Text--Cell Alignment}

Beyond aggregate metrics, per-type analysis reveals which cell types benefit most from text conditioning and where failures concentrate. Figure~\ref{fig:per-type-gen} presents per-type generation quality. The enhanced per-type analysis now makes heterogeneity explicit: centroid cosine is ranked across all 69 deduplicated cell types; the Fr\'echet profile highlights outlier types rather than only reporting the mean, and the abundance--fidelity scatter encodes Fr\'echet distance as bubble size and diversity ratio as color. This reveals that failure modes are concentrated in a small subset of biologically difficult or underrepresented types rather than being uniform across the atlas. Figure~\ref{fig:text-cell-alignment} shows the text--cell alignment analysis, confirming that generated cells remain most similar to their intended text condition despite these type-specific differences.

\begin{figure}[!htbp]
\centering
\includegraphics[width=0.95\textwidth]{figures/fig06_per_type_fidelity.pdf}
\caption{Per-type generation fidelity. (\textbf{a})~Per-type centroid cosine similarity ranked across cell types, (\textbf{b})~Fr\'{e}chet outlier profile with the global mean marked, and (\textbf{c})~abundance--fidelity scatter (bubble size $\propto$ Fr\'{e}chet distance, color encodes diversity ratio). Cell-type names are mapped by rank position in the supplementary materials.\label{fig:per-type-gen}}
\end{figure}

\vspace{-0.5em}
\begin{figure}[!htbp]
\centering
\includegraphics[width=0.95\textwidth]{figures/fig07_text_cell_alignment.pdf}
\caption{Text--cell alignment. (\textbf{a})~69$\times$69 cosine similarity heatmap between text and cell centroids, sorted by diagonal similarity. Mean diagonal similarity $\mu_{\text{diag}} = 0.916$ ($\sigma = 0.053$), mean off-diagonal $\mu_{\text{off}} = 0.044$ ($\sigma = 0.205$); Cohen's $d = 5.8$, separation ratio $20.9\times$, Mann--Whitney $U$ $p < 10^{-10}$. (\textbf{b})~Per-type alignment bars color-coded by quality tier (41 excellent $\geq 0.9$, 28 good, 0 poor). (\textbf{c})~Diagonal vs.\ off-diagonal distribution comparison with KDE overlay; the two populations are well separated ($\Delta = 0.872$, bootstrap 95\% CI [0.875, 0.913]).\label{fig:text-cell-alignment}}
\end{figure}

\subsection{Marker Gene Fidelity}

A key question for conditional generation is whether cell-type-specific marker gene signatures are preserved. Figure~\ref{fig:markers} compares the expression of canonical marker genes between real and generated cells, demonstrating that CLOP-DiT preserves cell-type-specific expression signatures. This directly shows that the biological signal resides in the intended marker genes rather than only in latent-space summary statistics.

\begin{figure}[!htbp]
\centering
\includegraphics[width=0.95\textwidth]{figures/fig08_marker_genes.pdf}
\caption{Marker gene comparison. (\textbf{a})~Paired bar chart of canonical marker gene expression in real (solid) vs.\ generated (hatched) cells, with bar colors grouping markers by lineage family (CD8 T, myeloid, epithelial, stromal); fold-change annotations flag markers deviating from unity. (\textbf{b})~Per-type $\times$ marker dual heatmap showing real and generated expression side by side for selected cell types. (\textbf{c})~Difference heatmap ($\Delta$ = generated $-$ real) with cell annotations for large deviations. (\textbf{d})~Fold-change waterfall for all markers, with a shaded band marking the $\pm$5\% agreement zone. Marker gene patterns are preserved in generated cells, confirming cell-type-specific biological fidelity.\label{fig:markers}}
\end{figure}

\subsection{Expression Correlation and Dimension Analysis}

To assess gene-level fidelity beyond marker genes, Figures~\ref{fig:expr-corr} and~\ref{fig:expr-analysis} quantify the per-gene Pearson correlation between real and generated mean expression across cell types and provide a broader view of expression distribution properties, including norm matching (generated mean norm 20.97 vs.\ real 20.89) and per-dimension statistics. Together they localize biological meaning in gene-wise means, variances, and marker-level residuals rather than only in embedding geometry.

\begin{figure}[!htbp]
\centering
\includegraphics[width=0.95\textwidth]{figures/fig09_expression_correlation.pdf}
\caption{Expression correlation analysis. (\textbf{a})~Per-gene mean expression scatter (real vs.\ generated), with point color encoding absolute residual magnitude and outlier genes annotated. (\textbf{b})~Per-type Pearson~$r$ lollipop chart with threshold-band shading. (\textbf{c})~Marker-gene expression comparison with error whiskers. (\textbf{d})~Per-gene residual distribution (generated $-$ real) with percentages within tolerance bands.\label{fig:expr-corr}}
\end{figure}

\vspace{-0.5em}
\begin{figure}[!htbp]
\centering
\includegraphics[width=0.95\textwidth]{figures/fig10_expression_analysis.pdf}
\caption{Expression distribution analysis. (\textbf{a})~Per-gene coefficient-of-variation scatter (real vs.\ generated) with annotated outliers. (\textbf{b})~Expression-range ribbons for real and generated cells. (\textbf{c})~Per-cell variability distributions. (\textbf{d})~Top variable genes ranked by standard-deviation ratio (generated/real), with a shaded unity band.\label{fig:expr-analysis}}
\end{figure}

\subsection{Conditioning Landscape and Diversity Trade-Offs}

Figure~\ref{fig:cond-umap} visualizes how text conditions distribute in the CLOP projection space and how the resulting generated cell clouds cluster accordingly. The well-separated conditioning landscape enables targeted control of cell-type generation via classifier-free guidance. Together with Figures~\ref{fig:diversity},~\ref{fig:noise-tradeoff}, and~\ref{fig:expr-diversity}, this section is the main robustness analysis for whether control can be increased without erasing biologically meaningful heterogeneity.

\begin{figure}[!htbp]
\centering
\includegraphics[width=0.95\textwidth]{figures/fig11_conditioning_umap.pdf}
\caption{Conditioning landscape. (\textbf{a--d})~PCA projections of the CLOP condition space for real cells and three conditioning modes (centroid, noisy, variant), showing cell counts and mean within-type diversity per mode. (\textbf{e})~Mean centroid shift from each generated mode to the real reference. (\textbf{f})~Within-type diversity by conditioning mode (dashed line = real baseline). (\textbf{g})~Per-type centroid shift distributions (box plots). (\textbf{h})~KNN-5 identity-preservation accuracy per conditioning mode. (\textbf{i})~Per-type diversity heatmap ($1 - \text{mean cosine}$). (\textbf{j})~Pairwise cosine violin plots comparing cluster tightness across modes.\label{fig:cond-umap}}
\end{figure}

Building on this conditioning structure, a sweep over eight CFG settings (0.5 to 5.0) with Euler and Midpoint ODE solvers reveals a clear accuracy--diversity trade-off (Figure~\ref{fig:diversity}). KNN accuracy saturates near CFG = 2.0--3.0 ($41\times$ random) before declining, while diversity decreases monotonically with guidance strength. The Midpoint solver at CFG = 1.0 achieves near-ideal diversity (0.930) while maintaining 80.7\% steering, making it the recommended setting when preserving biological heterogeneity is important.

\begin{figure}[!htbp]
\centering
\includegraphics[width=0.95\textwidth]{figures/fig12_diversity_diagnostics.pdf}
\caption{Diversity diagnostics. (\textbf{a})~Intra-type diversity ratio ranked across cell types (labels abbreviated for readability). (\textbf{b})~Nearest-neighbor distance distributions for real and generated populations. (\textbf{c})~CFG scale vs.\ diversity and norm trade-off across Euler and Midpoint solvers. (\textbf{d})~Centroid vs.\ noisy conditioning comparison. KNN accuracy saturates at CFG $\approx$ 2.0--3.0 while diversity decreases monotonically.\label{fig:diversity}}
\end{figure}

Figure~\ref{fig:noise-tradeoff} and Figure~\ref{fig:expr-diversity} further quantify the two operating regimes: high-fidelity (CFG $\geq$ 2.0, DivR $\approx$ 0.5) vs.\ high-diversity (CFG = 1.0 Midpoint, DivR = 0.93), and augment the global trade-off curves with per-type diversity-tail diagnostics. Expression-level diversity measurements confirm that the Midpoint solver best preserves gene-level variance.

\begin{figure}[!htbp]
\centering
\includegraphics[width=0.9\textwidth]{figures/fig13_noise_tradeoff.pdf}
\caption{Noise--scale trade-off. (\textbf{a})~Fr\'{e}chet distance (left axis), centroid cosine similarity, and diversity ratio (right axis) as a function of noise scale $\varepsilon$ at CFG = 1.5. The production configuration ($\varepsilon = 0.03$) is highlighted; a shaded band marks the sweet-spot region.\label{fig:noise-tradeoff}}
\end{figure}

\vspace{-0.5em}
\begin{figure}[!htbp]
\centering
\includegraphics[width=0.95\textwidth]{figures/fig14_expression_diversity.pdf}
\caption{Expression-level diversity. (\textbf{a})~Gene-expression variability summary comparing real and generated populations across key statistics. (\textbf{b})~Per-type gene standard-deviation ratio (generated/real), confirming that the recommended Midpoint solver preserves biologically meaningful variance.\label{fig:expr-diversity}}
\end{figure}

\subsection{Baseline Comparison and Composite Benchmark}
\vspace{-0.5em}
\FloatBarrier
\vspace{-0.5em}

CLOP-DiT is compared against baseline methods including unconditional generation (CFG = 0), Gaussian sampling from per-type statistics, shuffled-label and mean-collapse controls, EmbeddingVAE, scVI Latent~\cite{ref-scvi}, and a conditional scGen baseline~\cite{ref-scgen} (conditional scVI with cell-type covariates, trained on the same expression matrix; see Section~\ref{sec:discussion}). Figure~\ref{fig:baselines} shows that CLOP-DiT is strongest on conditioning-sensitive and downstream-relevant metrics (e.g., steering and classifier transfer), while some unconditional distribution-matching baselines may outperform on FD-like mean-matching objectives (see Discussion for the Fr\'{e}chet distance caveat).

\begin{figure}[!htbp]
\centering
\includegraphics[width=0.95\textwidth]{figures/fig15_baseline_comparison.pdf}
\caption{Baseline comparison. (\textbf{a})~Head-to-head grouped-bar analysis of CLOP-DiT vs.\ baselines (unconditional CFG\,=\,0, Gaussian per-type sampling, decoder-only) across four evaluation metrics (Fr\'{e}chet distance, centroid cosine similarity, diversity ratio, coverage). (\textbf{b})~Normalized radar chart showing multi-metric profiles for each method. (\textbf{c})~Relative improvement bar chart of CLOP-DiT over each baseline (values exceeding 300\% are capped for readability). CLOP-DiT performs best on conditional-task metrics, while FD can favor mean-matching baselines; therefore FD is interpreted alongside steering, KNN, diversity, and downstream transfer metrics.\label{fig:baselines}}
\end{figure}

Figure~\ref{fig:benchmark} aggregates these metrics into a composite benchmark score across eight methods (CLOP-DiT, scGen, scVI Latent, EmbeddingVAE, and four synthetic baselines). \textbf{On the 9 distributional metrics shared by all methods (the primary, fair like-for-like comparison), the Gaussian baseline (0.747) outperforms CLOP-DiT (0.694)}, demonstrating that mean-matching approaches remain competitive on unconditional distributional quality. The scGen baseline (common-metrics composite 0.425; trained as a conditional scVI with cell-type labels, equivalent to scGen's core capability) achieves the best \emph{coverage} among all methods (0.384 vs.\ 0.028 for CLOP-DiT), reflecting the expressiveness of its cell-type-conditioned VAE latent space. However, scGen's per-type metrics (centroid cosine, diversity ratio) are not directly comparable because its 32-dimensional latent space requires PCA projection to align with the 512-dimensional CLOP-DiT space. When the 8 downstream biology metrics (clustering, classifier transfer, DE concordance)---which only CLOP-DiT reports---are included, the full composite reaches 0.838 vs.\ 0.395 for Gaussian and 0.225 for scGen; however, this full composite is structurally biased in favour of CLOP-DiT because the 8 exclusive metrics guarantee a scoring advantage, and \textbf{should not be used for method ranking}. The common-metrics-only composite is the recommended comparison. Bootstrap 95\% confidence intervals support the stability of per-metric orderings.

\begin{figure}[!htbp]
\centering
\includegraphics[width=0.95\textwidth]{figures/fig16_benchmark.pdf}
\caption{Composite benchmark across eight methods including scGen and scVI Latent external baselines. (\textbf{a})~Normalized metrics heatmap (methods $\times$ metrics; green = high performance). (\textbf{b})~Composite score, with hatched bars denoting the common-metrics-only composite and solid bars denoting the full 17-metric composite. \textbf{Warning: the full 17-metric composite is structurally biased in favour of CLOP-DiT because 8 of 17 metrics are reported only by CLOP-DiT; the common-metrics-only (hatched) bars are the fair comparison.} (\textbf{c})~Grouped comparison of four interpretable benchmark metrics across methods. (\textbf{d})~95\% confidence intervals for FD, centroid cosine, and diversity ratio.\label{fig:benchmark}}
\end{figure}

\subsection{Downstream Biological Validation}
\vspace{-0.5em}
\FloatBarrier
\vspace{-0.5em}

Beyond distributional metrics, we evaluate whether generated cells are usable in standard single-cell analysis workflows. These analyses are implemented from matched real/generated AnnData objects with shared preprocessing in {\small\texttt{src.evaluation.\allowbreak{}downstream\_biology}} and use Scanpy~\cite{ref-scanpy} for the single-cell workflow, ensuring that any downstream advantage is not caused by mismatched feature engineering. Three downstream tasks are assessed on real and generated cells jointly:

\textbf{Clustering, mixing, and classifier alignment} (Figure~\ref{fig:downstream-pq}): We construct matched AnnData objects from real and generated expression matrices, select 2,000 shared highly variable genes, perform PCA and Leiden~\cite{ref-leiden} clustering on the combined dataset, and compute the Adjusted Rand Index (ARI), Normalized Mutual Information (NMI), cluster purity, and per-type $k$-NN mixing scores~\cite{ref-integration-benchmark}. High mixing scores indicate that generated cells partially co-localize with real cells of the same type in the joint embedding, though the discriminator AUC of 0.656 confirms that real and generated populations remain distinguishable. A logistic regression classifier is trained on real PCA embeddings to predict cell type, then evaluated on generated embeddings to measure transfer accuracy and macro-F1. The enhanced classifier panel reports a per-type precision/recall/F1 heatmap in addition to the confusion matrix and discriminator ROC, demonstrating that downstream performance exhibits strong heterogeneity across cell types. The real-vs.-generated discriminator AUC (0.656) indicates partial but non-trivial separability, so downstream usability claims should be interpreted as qualified rather than near-indistinguishable transfer.

\textbf{DE concordance} (Figure~\ref{fig:downstream-r}): Wilcoxon rank-sum differential expression analysis~\cite{ref-wilcoxon} is performed on three biologically meaningful contrasts (CD8 vs.\ CD4 T cells, resident macrophage vs.\ monocyte, epithelial tumor vs.\ fibroblast). We compare log-fold changes between real and generated data via Pearson and Spearman correlations, Jaccard overlap of the top-50 DE genes, and sign-agreement fractions. The enhanced DE figure weights each gene by effect size and colors it by adjusted significance, which separates minor noisy disagreements from the biologically important genes that dominate each contrast.

\begin{figure}[!htbp]
\centering
\includegraphics[width=0.95\textwidth]{figures/fig17_downstream_pq.pdf}
\caption{Downstream validation: clustering and classifier alignment. (\textbf{a})~Leiden clustering on combined real + generated data with UMAP overlay, (\textbf{b})~sorted per-type $k$-NN mixing profile (a supplementary table maps rank positions to cell-type names), and (\textbf{c})~compact cluster-alignment gauges (ARI, NMI, cluster purity, mean $k$-NN mixing). (\textbf{d})~Cell-type classification confusion matrix (real-trained on generated), (\textbf{e})~precision/recall/F1 heatmap restricted to the worst and best 20 cell types by F1, and (\textbf{f})~real-vs.-generated discriminator ROC curve. Median per-type F1 is reported in the heatmap stat box.\label{fig:downstream-pq}}
\end{figure}

\vspace{-0.5em}
\begin{figure}[!htbp]
\centering
\includegraphics[width=0.95\textwidth]{figures/fig18_de_concordance.pdf}
\caption{DE concordance. (\textbf{a})~Effect-size-weighted real-vs.-generated log-fold-change scatter for the leading contrast, with point size proportional to average absolute effect size and color indicating adjusted significance; sign-disagreement genes are outlined. (\textbf{b})~Compact summary heatmap reporting four metrics---Pearson~$r$, Spearman~$\rho$, Jaccard overlap of the top-50 DE genes (Jaccard@50), and sign agreement---across the three biologically meaningful contrasts (CD8 vs.\ CD4 T, resident macrophage vs.\ monocyte, epithelial vs.\ fibroblast). (\textbf{c})~Grouped-bar panel visualising the same per-contrast scores to aid comparison.\label{fig:downstream-r}}
\end{figure}
Together, these results establish that CLOP-DiT generates text-conditioned single-cell profiles with measurable type specificity, though with performance trade-offs between fidelity and diversity that depend on the operating regime.

\subsection{In-Distribution Expression-Level Validation}\label{sec:cross-dataset}

To assess biological fidelity across tissue contexts, we performed expression-level validation on five datasets spanning lung adenocarcinoma (GSE123902), basal cell carcinoma (GSE123813), liver T-cell hepatocellular carcinoma (GSE98638), gastric cancer (GSE183904), and acute lymphoblastic leukemia in bone marrow (GSE132509). Four of these five datasets (GSE123902, GSE123813, GSE98638, GSE132509) are from the training split; only GSE183904 is among the 8 held-out validation studies (Table~S2). This analysis therefore primarily evaluates in-distribution gene-level reconstruction fidelity rather than out-of-distribution generalization. For each dataset, cells were generated conditioned on the dataset-specific text descriptions and compared with real (scGPT-reconstructed) expression profiles; because both the generated and reference expressions are decoded by the same frozen scGPT decoder, the resulting correlations measure decoder-reconstructed fidelity rather than agreement with raw sequencing counts. Per-gene mean expression, gene-level variance, and Fr\'{e}chet distance in both gene space and embedding space were computed for each dataset independently (Table~\ref{tab:cross-dataset}).

\begin{table}[!htbp]
\caption{Cross-dataset biological validation (decoder-reconstructed fidelity). Per-gene mean expression Pearson correlation ($r$) and coefficient of determination ($R^2$), gene variance Pearson correlation ($r_\text{var}$), and Fr\'{e}chet distance in PCA-50 gene space (FD\textsubscript{gene}) and scGPT embedding space (FD\textsubscript{embed}) for real and generated cells. Both ``real'' and generated expression profiles are decoded by the same frozen scGPT decoder; correlations therefore measure decoder-reconstruction consistency rather than agreement with raw sequencing counts. Split indicates whether the dataset was used for training (Train) or held-out validation (Val). All five datasets achieve gene\_mean $r > 0.999$.\label{tab:cross-dataset}}
\centering
\small
\begin{tabular}{lccccccc}
\toprule
\textbf{Dataset (GEO)} & \textbf{Tissue} & \textbf{Split} & \textbf{$r_\text{mean}$} & \textbf{$R^2$} & \textbf{$r_\text{var}$} & \textbf{FD\textsubscript{gene}} & \textbf{FD\textsubscript{embed}} \\
\midrule
GSE123902 & Lung     & Train & 0.9995 & 0.9990 & $+$0.011  & $3.07 \times 10^{-5}$ & 73.9 \\
GSE123813 & Skin (BCC)  & Train & 0.9997 & 0.9995 & $-$0.019  & $5.25 \times 10^{-6}$ & 72.7 \\
GSE98638  & Liver    & Train & 0.9995 & 0.9989 & $-$0.003  & $2.04 \times 10^{-5}$ & 155.4 \\
GSE183904 & Gastric  & Val & 0.9997 & 0.9994 & $+$0.013  & $1.30 \times 10^{-5}$ & 62.1 \\
GSE132509 & Blood (ALL)  & Train & 0.9993 & 0.9985 & $-$0.007  & $8.87 \times 10^{-6}$ & 69.1 \\
\bottomrule
\end{tabular}
\end{table}

Across all five datasets, per-gene mean expression Pearson correlation exceeds 0.999 (range: 0.9993--0.9997), and $R^2$ exceeds 0.998 (range: 0.9985--0.9995). Because both reference and generated profiles pass through the same frozen scGPT decoder, these correlations reflect decoder-reconstructed fidelity: any latent vector near the training manifold decodes to a similar expression profile, so the high $r$ quantifies round-trip decoder consistency rather than absolute biological accuracy against raw counts. Fr\'{e}chet distance in PCA-50 gene space is consistently below $10^{-4}$ (range: $5.25 \times 10^{-6}$ to $3.07 \times 10^{-5}$), indicating excellent gene-space distributional agreement at the level of first and second moments.

However, per-gene variance correlation ($r_\text{var}$) is near zero or slightly negative across all datasets (range: $-0.019$ to $+0.013$), indicating that while the model reproduces mean expression accurately, it does not faithfully preserve gene-level variability structure. This is consistent with the diversity ratio analysis (Section~\ref{sec:results}): the flow matching objective optimizes for mean trajectory matching, and gene-level variance preservation may require explicit variance-matching terms in the loss function. The embedding-space Fr\'{e}chet distance (FD\textsubscript{embed}) ranges from 62.1 to 155.4 across datasets, substantially higher than gene-space FD. This apparent discrepancy---high decoder-reconstructed gene-mean fidelity alongside high embedding-space FD---arises because the scGPT decoder maps a broad range of latent vectors to similar expression profiles (many-to-one mapping), so gene-level fidelity does not require precise latent-space distributional matching.

\subsection{CLOP Ablation Study}\label{sec:ablation}

To validate the contribution of individual CLOP components, we conducted 15 ablation experiments, removing or modifying one component at a time while holding all other settings constant. Table~\ref{tab:ablation} summarizes the six most informative ablations.

\begin{table}[!htbp]
\caption{CLOP ablation study (\textbf{v8.2 baseline only}; not the production v9.3 model). Validation prototype accuracy (top-1) for the baseline configuration (v8.2, all features enabled) and representative ablations. $\Delta$~denotes the change from baseline in percentage points. These results characterize the CLOP alignment stage in isolation and do not directly predict downstream DiT generation quality; the production model (v9.3) uses different regularization settings.\label{tab:ablation}}
\centering
\small
\begin{tabular}{llcc}
\toprule
\textbf{Ablation} & \textbf{Category} & \textbf{Val Proto Acc (\%)} & \textbf{$\Delta$ (pp)} \\
\midrule
$-$cohesion        & regularization & 86.41 & $+$10.24 \\
$-$cell noise      & regularization & 86.23 & $+$10.07 \\
$-$dropout (lighter) & regularization & 77.32 & $+$1.15 \\
Baseline (v8.2)    & reference      & 76.17 & --- \\
Fixed temperature   & optimization   & 71.55 & $-$4.61 \\
High variant prob   & regularization & 70.14 & $-$6.03 \\
\bottomrule
\end{tabular}
\end{table}

The two most impactful ablations are cohesion regularization removal ($+$10.24~pp) and cell noise augmentation removal ($+$10.07~pp), both of which improve prototype accuracy. This counterintuitive result has a plausible mechanistic explanation: both cohesion loss and cell noise augmentation act as within-group tightening forces that penalize intra-cluster variance. In a contrastive learning regime where the primary bottleneck is inter-group separation (raw pairwise cosine 0.994), these tightening forces compete with the alignment gradient rather than complementing it. In effect, the network spends representational capacity enforcing tight within-group clusters before inter-group margins have been established, diverting training signal from the discrimination task. Removing these terms allows the optimizer to focus entirely on inter-group separation, which is the binding constraint at this scale. This interpretation is consistent with recent analyses of regularization--discrimination trade-offs in contrastive learning~\\cite{ref-siglip}. Architecture ablations (projection dimension 256 vs.\ 512 vs.\ 768) have minimal impact (within $\pm$0.3~pp), indicating that the CLOP aligner is not sensitive to projection dimension in this range. Fixed temperature ($-$4.61~pp) and high variant probability ($-$6.03~pp) degrade performance substantially relative to the v8.2 baseline (which used a learnable temperature). \textbf{Clarification:} The ablation baseline (v8.2) used a learnable $\tau$, whereas the production model (v9.3) uses fixed $\tau = 14.0$ in combination with a stratified validation split (Section~\ref{sec:dataset}). The production configuration was selected based on the full-pipeline generation quality assessment, not solely on CLOP prototype accuracy; the ablation table characterizes the CLOP alignment stage in isolation and does not directly predict downstream DiT generation performance.

\subsection{Rare Cell Augmentation}\label{sec:rare-cell}

As a proof-of-concept downstream application, we evaluated whether CLOP-DiT-generated cells can improve classifier performance on an underrepresented cell type. We identified ``Cycling cells'' as the rarest type in a test dataset (below 5\% prevalence), downsampled the training set to exacerbate the class imbalance, and augmented with synthetic cells at 1$\times$, 5$\times$, and 10$\times$ ratios. The baseline classifier (trained on real cells only) achieves an overall macro F1 of 0.964 and a Cycling-cell-specific F1 of 0.828. After augmentation at 1$\times$, overall macro F1 remains 0.953 and Cycling-cell F1 is 0.786; at 5$\times$ and 10$\times$, macro F1 stabilizes at 0.947 with Cycling-cell F1 at 0.741. The augmented cells therefore do not improve the rare-type classifier on this dataset, but they also do not substantially degrade overall classification, indicating that the synthetic profiles are sufficiently realistic to be mixed into training data without catastrophic interference. A more targeted augmentation strategy---selecting the most confident generated cells or combining with real-cell oversampling---may be needed to achieve a net positive augmentation effect.

\subsection{Conditioning Field Ablation and Swap-Label Permutation Test}\label{sec:cond-ablation}

To assess whether the conditioning signal causally steers generation---rather than the model recapitulating training-set structure regardless of prompt content---we performed two targeted experiments.

\textbf{Conditioning field ablation.} We generated 50 cells per type across all 69 cell types (3{,}450 cells per variant) under five prompt variants: (i)~\emph{full}: the complete structured prompt; (ii)~\emph{no-markers}: marker gene fields removed; (iii)~\emph{shuffled-markers}: marker gene lists randomly permuted across types; (iv)~\emph{metadata-only}: only organism, tissue, and disease fields retained; (v)~\emph{external-markers}: canonical markers from PanglaoDB/CellMarker~2.0 substituted for the training markers (evaluated on 15 well-characterized types). Table~\ref{tab:cond-ablation} reports kNN accuracy, steering accuracy, and centroid cosine similarity for each variant.

\begin{table}[htbp]
\centering
\caption{Conditioning field ablation results (CFG = 2.0, Euler solver, 10 steps). KNN accuracy assessed against real-cell references; steering measured as cosine alignment to the target type centroid; centroid cosine is the mean cosine similarity between generated and real centroids. The external-markers variant was evaluated on a 15-type subset with consensus markers from PanglaoDB and CellMarker~2.0.\label{tab:cond-ablation}}
\begin{tabular}{lrrr}
\toprule
Prompt variant & KNN (\%) & Steering (\%) & Centroid cos. \\
\midrule
Full (all fields) & 2.87 & 99.8 & 0.050 \\
No markers & 2.26 & 98.0 & 0.035 \\
Shuffled markers & 2.81 & 97.9 & 0.035 \\
Metadata only & 1.51 & 62.4 & 0.025 \\
External markers (15 types) & 17.07 & --- & 0.702$^*$ \\
\bottomrule
\multicolumn{4}{l}{\footnotesize $^*$Cosine similarity between external-marker and full-prompt centroids.}
\end{tabular}
\end{table}

The results demonstrate that semantic content drives the conditioning rather than prompt length alone: removing only the marker field reduces centroid cosine from 0.050 to 0.035 (30\% relative drop), while retaining only metadata fields collapses steering from 99.8\% to 62.4\%. Substituting external markers from independent databases yields a high cosine similarity to the full-prompt variant (0.702), indicating that the model responds to marker-gene semantics rather than memorized training prompts.

\textbf{Swap-label permutation test.} To test whether generation follows the conditioning signal causally, we generated 15 cell-type pairs where the prompt for type $A$ was used to condition generation but evaluation was performed against both type $A$ (matched) and type $B$ (mismatched) real-cell centroids. In 15 of 15 pairs, the generated cells were closer to the centroid of the \emph{conditioning} type than the evaluation type (matched cosine similarity 0.032 vs.\ mismatched 0.016, gap = 0.016). Moreover, in all 15 pairs the mislabeled cells followed the conditioning signal rather than the evaluation label, confirming that the model generates latents based on prompt semantics rather than memorized label identity. While this test does not fully exclude encoder-level leakage, it provides causal evidence that the conditioning channel is functional and semantically interpretable.

%%%%%%%%%%%%%%%%%%%%%%%%%%%%%%%%%%%%%%%%%%
\section{Discussion}\label{sec:discussion}

 CLOP-DiT shows that structured biological metadata can drive non-trivial single-cell latent generation, but the method should be interpreted narrowly and precisely. The conditioning signal is a fixed five-field template rather than free-form language, and the model is best understood as a structured-metadata-to-latent sampler rather than a general biological language model. Within that scope, the central result is clear: conditioning produces above-chance, cell-type-specific latent generation and supports two useful operating regimes that trade off fidelity against diversity.

\textbf{Decoder bottleneck and interpretation.} The most important limitation is the frozen scGPT decoder. Because this mapping is many-to-one, expression-level similarity is dominated by decoder behavior rather than by exact latent matching. This is why decoder-reconstructed mean expression can be highly concordant with real data even when latent-space Fr\'{e}chet distance remains large. The practical implication is simple: latent-space metrics are the primary evidence for generation quality, whereas expression-space results are best treated as downstream compatibility checks.

\textbf{Variance and covariance deficits.} The strongest biological limitation is the loss of cell-to-cell heterogeneity. Across cross-dataset validations, per-gene variance correlation is near zero, and within-type gene--gene correlation structure is only weakly preserved. In other words, CLOP-DiT captures conditional mean structure much better than conditional second-order structure. This helps explain why the model can preserve marker-level identity and gene means while still failing on applications that depend on realistic within-population variation. A variance-aware objective such as Equation~\ref{eq:var-loss}, or a latent-space distributional penalty, is therefore a more compelling next step than simply scaling model size.

The standard flow matching loss (Equation~\ref{eq:fm}) optimizes $\| v_\theta(z_t, t, c) - (z_1 - z_0) \|^2$, which regresses toward the conditional mean velocity. While individual trajectories $z_0 \to z_1$ are stochastic, the MSE objective rewards accurate prediction of the expected velocity at each $(z_t, t, c)$ triplet without penalising insufficient spread in the generated distribution. An explicit \textbf{variance-matching regularisation} term of the form
\begin{linenomath}
\begin{equation}
  \mathcal{L}_\text{var} = \sum_{g=1}^{G} \left( \sigma_g^{\text{gen}} - \sigma_g^{\text{real}} \right)^2,
  \label{eq:var-loss}
\end{equation}
\end{linenomath}
where $\sigma_g^{\text{gen}}$ and $\sigma_g^{\text{real}}$ denote the per-gene standard deviations of generated and real cells respectively, could be added as a post-hoc penalty on mini-batch statistics during DiT training. In practice this would require differentiable decoding within the training loop---currently infeasible with the frozen scGPT decoder---or, more realistically, a latent-space distributional penalty such as the sliced Wasserstein distance between real and generated distributions for each cell type.

\textbf{Pilot variance-matching analysis.} A preliminary latent-space analysis using sliced Wasserstein distance (SWD; Figure~\ref{fig:var-pilot}) supports this interpretation: generated latents preserve overall spread better than per-dimension variance structure. That makes latent-space regularisation attractive because it targets the generator directly instead of asking the frozen decoder to adjudicate generative quality. Implementing such a penalty is left to future work.

\begin{figure}[!htbp]
\centering
\includegraphics[width=0.95\textwidth]{figures/fig19_variance_matching_pilot.pdf}
\caption{Pilot variance-matching analysis. (\textbf{a})~Per-cell-type sliced Wasserstein distance (SWD) between real and generated latent distributions. (\textbf{b})~Per-type latent variance ratios (generated/real). (\textbf{c})~Per-dimension variance correlations. (\textbf{d})~SWD versus training cell count. The figure shows partial preservation of variance magnitude without strong structural alignment.\label{fig:var-pilot}}
\end{figure}

\textbf{Gene--gene structure.} Figure~\ref{fig:gene-gene-corr} extends the same conclusion to covariance: within-type gene--gene correlation matrices are only weakly preserved. This is consistent with the variance results rather than a separate failure mode. Both point to the same limitation, namely that the current objective matches central tendency far better than biologically meaningful heterogeneity.

\begin{figure}[!htbp]
\centering
\includegraphics[width=0.95\textwidth]{figures/fig20_gene_gene_correlation.pdf}
\caption{Gene--gene correlation preservation across cell types. (\textbf{a})~Distribution of per-type Mantel correlations between real and generated correlation matrices. (\textbf{b--c})~Correlation-difference heatmaps for representative best- and worst-preserved cell types. (\textbf{d})~Mantel $r$ versus real-cell count per type. The weak Mantel correlations indicate limited preservation of within-type co-expression structure.\label{fig:gene-gene-corr}}
\end{figure}

\textbf{Benchmarking and operating regimes.} CLOP-DiT does not uniformly dominate the benchmark. On the primary common-metrics-only comparison, the Gaussian baseline outperforms CLOP-DiT, which underscores that mean-matching remains competitive when evaluation ignores conditioning-sensitive biology. The more informative result is that CLOP-DiT creates a controllable trade-off: CFG = 2.0 with Euler sampling yields stronger type fidelity, whereas CFG = 1.0 with Midpoint sampling preserves substantially more diversity. This trade-off is useful, but it does not erase the remaining gap to real data, as reflected by the 36.9\% KNN accuracy, the 52-point gap from the real-data reference, and the discriminator AUC of 0.656.

\textbf{Practical impact of KNN misclassification.} The 36.9\% KNN accuracy, while $25\times$ above random chance, means that roughly two-thirds of generated cells are assigned to the wrong cell type by a nearest-neighbor classifier. Concretely, a downstream user who treats generated cells as labeled training data for a supervised classifier would inherit an approximate 63\% label-noise rate in latent space. For applications requiring high type-label confidence---for example, augmenting training sets for rare-cell classifiers or seeding perturbation-response simulations---the current accuracy is insufficient without post-hoc filtering (e.g., retaining only the top-confidence cells by cosine proximity to the target centroid). Improving KNN accuracy is therefore a prerequisite for most practical single-cell generation pipelines.

\textbf{Biological consequence of diversity loss.} The high-fidelity regime (CFG = 2.0) achieves a diversity ratio of 0.513, meaning generated cells retain roughly half the within-type variance of real data. Biologically, this halved variance implies that the generated distribution collapses toward cell-type centroids, erasing the within-population heterogeneity that encodes functional cell states (e.g., activation, exhaustion, and cycling gradients within T-cell populations). Any downstream analysis relying on within-type substructure---trajectory inference, gene regulatory network inference, or sub-clustering into functional states---would be unreliable on data generated in the high-fidelity regime. The high-diversity regime (DivR = 0.93) largely mitigates this issue but at the cost of reduced type specificity (KNN $\approx$ 29\%).

\textbf{Scope and unresolved limitations.}\label{sec:limitations} The current evidence is restricted to human and mouse cancer/developmental datasets and does not constitute a stringent out-of-distribution test. Direct comparison with text-centric systems such as Cell2Sentence remains future work because preprocessing and evaluation pipelines differ substantially. External marker validation supports the biological plausibility of the conditioning markers, but because markers are derived from the training corpus and the decoder masks prompt-level expression differences, residual conditioning leakage cannot be ruled out completely.

\textbf{Negative result: rare-cell augmentation.} The pilot rare-cell augmentation experiment is important rather than incidental: generated cells did not improve rare-type classification. This failure is consistent with the variance results above. If generated cells behave mostly as mean-expression representatives, then augmentation adds centroid-like examples without the within-type heterogeneity needed to broaden the training distribution. Practical augmentation will therefore require both better variance preservation and a decoder that maps distinct latent states to genuinely distinct expression profiles.

%%%%%%%%%%%%%%%%%%%%%%%%%%%%%%%%%%%%%%%%%%
\section{Conclusions}\label{sec:conclusions}

We introduced CLOP-DiT, a two-stage pipeline for structured-metadata-conditioned single-cell latent generation. CLOP aligns BiomedBERT text embeddings and scGPT cell embeddings in a shared 512-dimensional space, and a conditional Diffusion Transformer then samples scGPT latents from a fixed five-field biological template. The method is therefore best interpreted as a structured-metadata-to-latent conditional sampler rather than a free-form language model. Across 69 deduplicated cell types from 80 datasets, CLOP-DiT achieves 36.9\% KNN accuracy and 81.0\% steering in a high-fidelity regime, while a high-diversity regime reaches a diversity ratio of 0.93 at 80.7\% steering. Multi-seed replication confirms that steering and diversity are stable across independent runs. A conditioning field ablation and swap-label permutation test provide causal evidence that the conditioning signal is semantically functional.

\textbf{Scope of claims.} All results are limited to in-distribution human and mouse cancer and developmental datasets; no stringent out-of-distribution generalization test has been conducted. The evaluation is restricted to the 69 deduplicated cell types and 80 GEO datasets described in Section~\ref{sec:dataset}. Claims about generative quality refer exclusively to latent-space metrics; expression-level concordance is dominated by the frozen scGPT decoder and should not be interpreted as evidence of faithful transcriptomic generation.

\textbf{Future directions.} Concrete next steps include: (i)~\emph{variance-aware training}: incorporating a latent-space distributional penalty (e.g., sliced Wasserstein distance) into the DiT training objective to preserve within-type heterogeneity; (ii)~\emph{decoder replacement}: replacing the frozen scGPT decoder with a fine-tuned or jointly trained decoder that better preserves the diversity encoded in the latent space; (iii)~\emph{out-of-distribution evaluation}: testing conditioning on cell types, tissues, and organisms absent from the training set; (iv)~\emph{covariance-aware generation}: adding explicit gene--gene correlation matching to the training objective; (v)~\emph{downstream task validation}: evaluating the practical utility of generated cells for data augmentation, perturbation prediction, and trajectory inference on benchmarks with ground-truth labels.

%%%%%%%%%%%%%%%%%%%%%%%%%%%%%%%%%%%%%%%%%%

%%%%%%%%%%%%%%%%%%%%%%%%%%%%%%%%%%%%%%%%%%
%% optional
\supplementary{The following supporting information can be downloaded at: \linksupplementary{s1}. Table~S1: GEO accession identifiers for all 80 training datasets; Table~S2: held-out validation dataset identifiers; Table~S3: full CFG sweep results (8 guidance scales $\times$ 2 solvers $\times$ 3 metrics); Table~S4: biological validation datasets; Table~S5: bootstrap 95\% confidence intervals ($B{=}1{,}000$, percentile method) for headline generation metrics. Multi-seed validation data (3 independent seeds), independent marker gene validation data (PanglaoDB and CellMarker~2.0 cross-references), and external marker ablation analysis are available in the repository at \texttt{results/multi\_seed/} and \texttt{results/marker\_independence/}. Reproduction instructions (environment setup, \texttt{scripts/pipeline/regenerate\_report.sh}, expected outputs) are provided in \texttt{docs/operational/QUICK\_START.md} and \texttt{scripts/pipeline/regenerate\_report.sh}. All figures use colormaps designed for accessibility under common forms of color vision deficiency (deuteranopia, protanopia); critical distinctions additionally use shape or pattern encoding.}

% Only for journal Methods and Protocols:
% If you wish to submit a video article, please do so with any other supplementary material.
% \supplementary{The following supporting information can be downloaded at: \linksupplementary{s1}, Figure S1: title; Table S1: title; Video S1: title. A supporting video article is available at doi: link.}

% Only used for preprtints:
% \supplementary{The following supporting information can be downloaded at the website of this paper posted on \href{https://www.preprints.org/}{Preprints.org}.}

% Only for journal Hardware:
% If you wish to submit a video article, please do so with any other supplementary material.
% \supplementary{The following supporting information can be downloaded at: \linksupplementary{s1}, Figure S1: title; Table S1: title; Video S1: title.\vspace{6pt}\\
%\begin{tabularx}{\textwidth}{lll}
%\toprule
%\textbf{Name} & \textbf{Type} & \textbf{Description} \\
%\midrule
%S1 & Python script (.py) & Script of python source code used in XX \\
%S2 & Text (.txt) & Script of modelling code used to make Figure X \\
%S3 & Text (.txt) & Raw data from experiment X \\
%S4 & Video (.mp4) & Video demonstrating the hardware in use \\
%... & ... & ... \\
%\bottomrule
%\end{tabularx}
%}

%%%%%%%%%%%%%%%%%%%%%%%%%%%%%%%%%%%%%%%%%%
\authorcontributions{Z.F.: conceptualization, methodology, software, validation, formal analysis, investigation, resources, data curation, writing---original draft preparation, writing---review and editing, visualization. The author has read and agreed to the published version of the manuscript.}

\funding{This research received no external funding.}

\institutionalreview{Ethical review and approval were waived for this study, as only publicly available, de-identified data from the Gene Expression Omnibus (GEO) were analyzed; all datasets were deposited by their original investigators under institutional ethical approvals as required by GEO submission policy. No new human subjects research, animal experiments, or clinical interventions were conducted.}

\informedconsent{Not applicable.}

\dataavailability{The training and validation data were curated from publicly available GEO datasets; a full list of GEO accession identifiers is provided in Supplementary Table~S1, and the held-out validation study identifiers are in Supplementary Table~S2. The train/validation split configuration is defined in \texttt{configs/clop.yaml} and is deterministic given the dataset identifiers. Source code, trained model checkpoints (CLOP aligner and DiT generator), configuration files, and the figure-reproduction script (\texttt{scripts/pipeline/regenerate\_report.sh}) are publicly available on GitHub~\cite{ref-clop-dit} and will be archived on Zenodo upon acceptance of this manuscript. The complete software environment is specified in \texttt{requirements.txt} (Python~3.10, PyTorch~2.1, CUDA~12.0, scGPT~v0.2.1, Transformers~4.36). One-command figure reproduction from cached embeddings is documented in \texttt{docs/operational/QUICK\_START.md}.}

%\durcstatement{Current research is limited to the [please insert a specific academic field, e.g., XXX], which is beneficial [share benefits and/or primary use] and does not pose a threat to public health or national security. Authors acknowledge the dual-use potential of the research involving xxx and confirm that all necessary precautions have been taken to prevent potential misuse. As an ethical responsibility, authors strictly adhere to relevant national and international laws about DURC. Authors advocate for responsible deployment, ethical considerations, regulatory compliance, and transparent reporting to mitigate misuse risks and foster beneficial outcomes.}

% Only for journal Nursing Reports
%\publicinvolvement{Please describe how the public (patients, consumers, carers) were involved in the research. Consider reporting against the GRIPP2 (Guidance for Reporting Involvement of Patients and the Public) checklist. If the public were not involved in any aspect of the research add: ``No public involvement in any aspect of this research''.}
%
%% Only for journal Nursing Reports
%\guidelinesstandards{Please add a statement indicating which reporting guideline was used when drafting the report. For example, ``This manuscript was drafted against the XXX (the full name of reporting guidelines and citation) for XXX (type of research) research''. A complete list of reporting guidelines can be accessed via the equator network: \url{https://www.equator-network.org/}.}
%
%% Only for journal Nursing Reports
%\useofartificialintelligence{Please describe in detail any and all uses of artificial intelligence (AI) or AI-assisted tools used in the preparation of the manuscript. This may include, but is not limited to, language translation, language editing and grammar, or generating text. Alternatively, please state that “AI or AI-assisted tools were not used in drafting any aspect of this manuscript”.}

\acknowledgments{The author acknowledges the developers and communities behind scGPT, BiomedBERT, and the Gene Expression Omnibus (GEO) for providing open-access tools and datasets that made this work possible.}

\conflictsofinterest{The author declares no conflicts of interest.}

%%%%%%%%%%%%%%%%%%%%%%%%%%%%%%%%%%%%%%%%%%
%% Optional

%% Only for journal Encyclopedia
%\entrylink{The Link to this entry published on the encyclopedia platform.}

\abbreviations{Abbreviations}{
The following abbreviations are used in this manuscript:
\\

\noindent
\begin{tabular}{@{}ll}
AdaLN & Adaptive Layer Normalization\\
ALL & Acute Lymphoblastic Leukemia\\
ALS & Amyotrophic Lateral Sclerosis\\
AML & Acute Myeloid Leukemia\\
AMP & Automatic Mixed Precision\\
ARI & Adjusted Rand Index\\
ASD & Autism Spectrum Disorder\\
AUC & Area Under Curve\\
BCC & Basal Cell Carcinoma\\
BERT & Bidirectional Encoder Representations from Transformers\\
CE & Cross-Entropy\\
CFG & Classifier-free guidance\\
CI & Confidence Interval\\
CLIP & Contrastive Language-Image Pre-training\\
CLOP & Contrastive Language-Omics Pretraining\\
CRISPR & Clustered Regularly Interspaced Short Palindromic Repeats\\
CUDA & Compute Unified Device Architecture\\
DE & Differential expression\\
DiT & Diffusion Transformer\\
DivR & Diversity Ratio\\
EMA & Exponential Moving Average\\
ESC & Embryonic Stem Cell\\
FD & Fr\'{e}chet Distance\\
FP16 & Half-Precision Floating Point\\
GAN & Generative Adversarial Network\\
GELU & Gaussian Error Linear Unit\\
GEO & Gene Expression Omnibus\\
GPU & Graphics Processing Unit\\
HBV & Hepatitis B Virus\\
IFN & Interferon\\
InfoNCE & Info Noise-Contrastive Estimation\\
KNN & $k$-Nearest Neighbor\\
LinAcc & Linear Classifier Accuracy\\
logFC & Log Fold Change\\
LoRA & Low-Rank Adaptation\\
LPS & Lipopolysaccharide\\
LR & Learning Rate\\
MLP & Multi-Layer Perceptron\\
MSE & Mean Squared Error\\
NK & Natural Killer\\
NMI & Normalized Mutual Information\\
ODE & Ordinary Differential Equation\\
OOD & Out-of-Distribution\\
PBMC & Peripheral Blood Mononuclear Cell\\
PCA & Principal Component Analysis\\
ReLU & Rectified Linear Unit\\
ROC & Receiver Operating Characteristic\\
scGPT & Single-cell Generative Pre-trained Transformer\\
scRNA-seq & Single-cell RNA sequencing\\
SigLIP & Sigmoid Loss for Language-Image Pre-training\\
UMAP & Uniform Manifold Approximation and Projection\\
VAE & Variational Autoencoder\\
VRAM & Video Random Access Memory\\
ZCA & Zero-phase Component Analysis
\end{tabular}
}

%%%%%%%%%%%%%%%%%%%%%%%%%%%%%%%%%%%%%%%%%%
%% Optional
\appendixtitles{yes}
\appendixstart
\appendix
\section[\appendixname~\thesection]{Architecture Details}

Table~\ref{tab:architecture} summarizes the complete architecture specifications of the CLOP-DiT pipeline.

\begin{table}[htbp]
\caption{Architecture specifications.\label{tab:architecture}}
\begin{tabularx}{\textwidth}{lCC}
\toprule
\textbf{Component} & \textbf{Architecture} & \textbf{Parameters} \\
\midrule
BiomedBERT-large (frozen) & Transformer encoder  & 340M \\
CLOP text projector       & MLP [1024$\to$1024$\to$1024$\to$512] & 1.58M \\
CLOP cell projector       & MLP [512$\to$512$\to$512$\to$512]    & 1.44M \\
DiT1D                     & 8$\times$ AdaLN-Zero blocks, 8 heads & 22.10M \\
scGPT encoder (frozen)    & Transformer                          & $\sim$51M \\
scGPT decoder (frozen)    & Gene token prediction head           & included \\
\bottomrule
\end{tabularx}
\end{table}

\section[\appendixname~\thesection]{CFG Scale Sweep}

Table~\ref{tab:cfg-sweep} reports the full CFG sweep across 8 guidance scales with 10-step Euler integration and 69 evaluation groups. Note that these values are computed on a different subsample (69 groups $\times$ 200 cells) than the primary operating-point results in Table~\ref{tab:core-results} (69 types $\times$ 200 cells with different random seeds), which accounts for minor numerical differences between the two tables.

\begin{table}[htbp]
\caption{CFG scale sweep (10-step Euler, 69 eval groups).\label{tab:cfg-sweep}}
\centering
\begin{tabular}{lrrrr}
\toprule
\textbf{CFG} & \textbf{KNN-1} & \textbf{Steering} & \textbf{DivR} & \textbf{KNN/Random} \\
\midrule
0.5 & 0.147 & 0.788 & 1.13 & 15$\times$ \\
1.0 & 0.348 & 0.796 & 0.76 & 35$\times$ \\
1.5 & 0.396 & 0.798 & 0.60 & 40$\times$ \\
2.0 & 0.407 & 0.806 & 0.53 & 41$\times$ \\
2.5 & 0.408 & 0.804 & 0.49 & 41$\times$ \\
3.0 & 0.410 & 0.802 & 0.48 & 41$\times$ \\
4.0 & 0.402 & 0.804 & 0.50 & 40$\times$ \\
5.0 & 0.380 & 0.798 & 0.55 & 38$\times$ \\
\bottomrule
\end{tabular}
\end{table}

\section[\appendixname~\thesection]{Training Hyperparameters}\label{app:hyperparams}

Table~\ref{tab:hyperparams} lists the complete training hyperparameters for both CLOP alignment and DiT flow-matching stages, corresponding to \texttt{configs/clop.yaml} and \texttt{configs/dit.yaml}.

\begin{table}[htbp]
\caption{Training hyperparameters.\label{tab:hyperparams}}
\begin{tabularx}{\textwidth}{lCC}
\toprule
\textbf{Hyperparameter} & \textbf{CLOP} & \textbf{DiT} \\
\midrule
Optimizer         & AdamW              & AdamW ($\beta_1{=}0.9, \beta_2{=}0.999$) \\
Learning rate     & $5 \times 10^{-4}$ & $2 \times 10^{-4}$ \\
Weight decay      & 0.06               & 0.01 \\
Batch size        & 1024               & 1024 \\
Epochs            & 60                 & 200 \\
Warmup epochs     & 5                  & 5 \\
LR schedule       & Cosine with warmup & Cosine with warmup \\
Gradient clipping  & 1.0               & 1.0 \\
Mixed precision   & FP16 (AMP)         & FP16 (AMP) \\
Dropout           & 0.2                & Proj: 0.1, Attn: 0.0 \\
EMA decay         & ---                & 0.9999 \\
Temperature       & 14.0 (fixed; see Section~\ref{sec:clop})       & --- \\
Loss              & PrototypeSigLIP    & Flow matching MSE \\
Cohesion weight   & 0.15               & --- \\
Label smoothing   & 0.1                & --- \\
CFG drop prob     & ---                & 0.15 \\
Time sampling     & ---                & Logit-normal ($\mu{=}0, \sigma{=}1$) \\
\bottomrule
\end{tabularx}
\end{table}

\section[\appendixname~\thesection]{Supplementary Tables}\label{app:suppl-tables}

% Table S1: All 80 GEO Dataset Identifiers
{\footnotesize
\renewcommand{\arraystretch}{0.9}
\begin{longtable}{rlp{5.5cm}rr}
\caption{All 80 GEO datasets used in CLOP-DiT training and validation. Datasets were preprocessed to 2{,}000 highly variable genes per dataset. Split: Train or held-out Validation (study-level, seed 42).\label{tab:s1_datasets}} \\
\toprule
\# & Accession & Description & $n_{\text{cells}}$ & Split \\
\midrule
\endfirsthead
\toprule
\# & Accession & Description & $n_{\text{cells}}$ & Split \\
\midrule
\endhead
\midrule
\multicolumn{5}{r}{\textit{Continued on next page}} \\
\endfoot
\bottomrule
\endlastfoot
  1 & GSE103867 & scRNA-seq (H.~sapiens) & 2,787 & Val \\
  2 & GSE110686 & scRNA-seq (H.~sapiens) & 2,992 & Train \\
  3 & GSE115571 & Mouse cells, LPS stimulation & 442 & Train \\
  4 & GSE117988 & PBMCs from patients & 2,638 & Train \\
  5 & GSE117988 & Merkel cell carcinoma tumor & 2,990 & Train \\
  6 & GSE120505 & Aged blood cells & 3,000 & Train \\
  7 & GSE123813 & Human basal cell carcinoma & 3,000 & Train \\
  8 & GSE123813 & Human cutaneous squamous cell carcinoma & 3,000 & Train \\
  9 & GSE123902 & Human lung adenocarcinoma & 2,753 & Train \\
  10 & GSE124310 & Human multiple myeloma bone marrow & 2,833 & Train \\
  11 & GSE130148 & Human fetal lung development & 3,000 & Train \\
  12 & GSE132509 & Pediatric bone marrow mononuclear cells & 2,984 & Train \\
  13 & GSE138709 & Human intrahepatic cholangiocarcinoma & 2,891 & Train \\
  14 & GSE142653 & Human pituitary gland development & 2,885 & Train \\
  15 & GSE143423 & Human leptomeningeal brain metastasis & 3,000 & Train \\
  16 & GSE143423 & Triple-neg.\ breast cancer brain met. & 2,332 & Train \\
  17 & GSE145929 & Adult mouse prostate & 2,830 & Train \\
  18 & GSE145929 & Adult mouse urethra & 2,745 & Train \\
  19 & GSE148218 & Human bone marrow, ALL & 2,807 & Train \\
  20 & GSE149655 & Human early-stage lung adenocarcinoma & 2,087 & Train \\
  21 & GSE155109 & Endothelial cells, breast cancer & 3,000 & Train \\
  22 & GSE155109 & Stromal cells, breast cancer & 3,000 & Train \\
  23 & GSE163558 & Human stomach cancer & 2,467 & Train \\
  24 & GSE165784 & Human retina development & 2,622 & Train \\
  25 & GSE165844 & Mouse LSK hematopoietic progenitors & 2,881 & Train \\
  26 & GSE167597 & Mouse spinal cord development & 3,000 & Val \\
  27 & GSE168181 & Human breast cancer tissue & 2,863 & Train \\
  28 & GSE175975 & scRNA-seq (H.~sapiens) & 2,728 & Train \\
  29 & GSE183904 & Human gastric cancer & 3,000 & Val \\
  30 & GSE189070 & Mouse astrocytes, spinal cord injury & 2,960 & Train \\
  31 & GSE189357 & Human lung adenocarcinoma & 2,886 & Train \\
  32 & GSE192857 & hESC differentiation time-series & 2,696 & Train \\
  33 & GSE212502 & Human forebrain organoids & 3,000 & Train \\
  34 & GSE213740 & Human ascending aortic wall & 2,704 & Val \\
  35 & GSE220913 & Human HL-60 cell lines, CFTR & 2,972 & Train \\
  36 & GSE222002 & T cells from human cancer & 2,962 & Train \\
  37 & GSE222369 & NK cells, human lymphoma & 2,778 & Train \\
  38 & GSE225600 & Human breast cancer tumor & 2,333 & Train \\
  39 & GSE225857 & Colorectal cancer liver metastases & 2,986 & Train \\
  40 & GSE226131 & Aged mouse hematopoietic stem cells & 2,608 & Train \\
  41 & GSE226762 & scRNA-seq (H.~sapiens) & 830 & Val \\
  42 & GSE227644 & scRNA-seq (H.~sapiens) & 3,000 & Train \\
  43 & GSE227719 & Mouse KrasG12D/P53 lung organoids & 2,990 & Train \\
  44 & GSE228499 & Human breast cancer & 2,448 & Train \\
  45 & GSE235787 & B cells, ALL & 2,798 & Train \\
  46 & GSE236565 & Mouse splenic CD4+ T cells & 2,556 & Train \\
  47 & GSE248214 & scRNA-seq (H.~sapiens) & 2,998 & Train \\
  48 & GSE253355 & Human bone marrow niche & 2,835 & Train \\
  49 & GSE262288 & Metastatic breast cancer & 2,321 & Train \\
  50 & GSE272187 & scRNA-seq (H.~sapiens) & 2,639 & Train \\
  51 & GSE275119 & Mouse dental tissue development & 2,963 & Train \\
  52 & GSE277740 & Mouse thalamus, SBS vs DNBSEQ & 2,644 & Val \\
  53 & GSE279781 & scRNA-seq (H.~sapiens) & 2,654 & Train \\
  54 & GSE280847 & CD45+ leukocytes, mouse MC38 tumor & 2,864 & Train \\
  55 & GSE283205 & Hepatoblastoma tumor tissue & 2,994 & Train \\
  56 & GSE283397 & Human CD8+ T cells, chronic HBV & 2,856 & Train \\
  57 & GSE288211 & Mouse pancreatic islets, Zzef1 KO & 2,724 & Train \\
  58 & GSE289611 & Human pulmonary pleomorphic carcinoma & 2,983 & Train \\
  59 & GSE291166 & Mouse mesenteric lymph node & 2,901 & Train \\
  60 & GSE299623 & scRNA-seq (H.~sapiens) & 2,991 & Train \\
  61 & GSE300862 & Mouse Kras/Trp53 lung adenocarcinoma & 2,716 & Val \\
  62 & GSE302433 & Mouse B16F10 melanoma, ZDHHC13 & 2,953 & Val \\
  63 & GSE303309 & Brain CD45+ immune cells, Alzheimer's & 2,999 & Train \\
  64 & GSE306567 & Mouse ESCs and neural progenitors & 2,454 & Train \\
  65 & GSE306676 & Mouse spinal cord, SOD1-G93A ALS & 2,959 & Train \\
  66 & GSE307261 & Human B and T cells, melanoma & 2,921 & Train \\
  67 & GSE307774 & Mouse colorectal tumors, Kras/Braf & 2,505 & Train \\
  68 & GSE308428 & Mouse hippocampal microglia, ASD & 2,991 & Train \\
  69 & GSE309368 & Human myeloid cells, healthy donors & 1,361 & Train \\
  70 & GSE315628 & scRNA-seq (M.~musculus) & 2,895 & Train \\
  71 & GSE318353 & scRNA-seq (H.~sapiens) & 2,964 & Train \\
  72 & GSE98638 & Human hepatocellular carcinoma T cells & 3,000 & Train \\
  73 & GSE120446 & Bone marrow, hematopoietic development & 2,890 & Train \\
  74 & GSE104323 & Dentate gyrus, brain development & 3,000 & Train \\
  75 & GSE75748 & Endoderm differentiation & 2,531 & Train \\
  76 & GSE144024 & Human embryonic stem cells & 3,000 & Train \\
  77 & GSE72857 & Hematopoiesis & 3,000 & Train \\
  78 & GSE226824 & IFN-treated hematopoietic progenitors & 2,739 & Train \\
  79 & GSE141259 & Lung development & 3,000 & Train \\
  80 & GSE139369 & Human hematopoiesis (Setty et al.) & 2,995 & Train \\
\midrule
  & \textbf{Total} & & \textbf{220,304} & \\
\end{longtable}
}% end footnotesize for Table S1

% Table S2: 8 Held-out Validation Datasets
\begin{table}[htbp]
\centering
\caption{The 8 held-out validation datasets (study-level split, seed 42, val\_split = 0.1). No cells from these studies were seen during CLOP or DiT training.\label{tab:s2_validation}}
\begin{tabular}{rlp{5.5cm}r}
\toprule
\# & Accession & Description & $n_{\text{cells}}$ \\
\midrule
  1 & GSE103867 & scRNA-seq (H.~sapiens) & 2,787 \\
  2 & GSE167597 & Mouse spinal cord development & 3,000 \\
  3 & GSE183904 & Human gastric cancer & 3,000 \\
  4 & GSE213740 & Human ascending aortic wall & 2,704 \\
  5 & GSE226762 & scRNA-seq (H.~sapiens) & 830 \\
  6 & GSE277740 & Mouse thalamus, SBS vs DNBSEQ & 2,644 \\
  7 & GSE300862 & Mouse Kras/Trp53 lung adenocarcinoma & 2,716 \\
  8 & GSE302433 & Mouse B16F10 melanoma, ZDHHC13 & 2,953 \\
\midrule
  & \textbf{Total} & & \textbf{20,634} \\
\bottomrule
\end{tabular}
\end{table}

% Table S3: Full CFG Scale / Solver / Steps Sweep Results
{\footnotesize
\renewcommand{\arraystretch}{0.9}
\begin{longtable}{lllrrrrrr}
\caption{Full CFG scale sweep. 200 cells generated per group (69 groups). KNN: type-identity accuracy; Steer.: alignment; DivR: diversity ratio (1.0 = ideal); LinAcc: linear classifier accuracy; FD: Fr\'{e}chet distance.\label{tab:s3_cfg_sweep}} \\
\toprule
CFG & Solver & Steps & KNN-1 & KNN-5 & Steer. & DivR & LinAcc & FD \\
\midrule
\endfirsthead
\toprule
CFG & Solver & Steps & KNN-1 & KNN-5 & Steer. & DivR & LinAcc & FD \\
\midrule
\endhead
\midrule
\multicolumn{9}{r}{\textit{Continued on next page}} \\
\endfoot
\bottomrule
\endlastfoot
\multicolumn{9}{l}{\textit{Primary configurations (Table~\ref{tab:core-results})}} \\
\midrule
  3.0 & Euler & 10 & 0.367 & 0.554 & 0.802 & 0.473 & 0.508 & 3.46 \\
  2.0 & Euler & 10 & 0.369 & 0.553 & 0.810 & 0.513 & 0.511 & 3.17 \\
  1.0 & Euler & 10 & 0.313 & 0.483 & 0.810 & 0.745 & 0.390 & 2.84 \\
  0.0 & Euler & 10 & 0.010 & 0.051 & 0.475 & 1.833 & 0.010 & 0.58 \\
  3.0 & Midpoint & 10 & 0.340 & 0.528 & 0.792 & 0.628 & 0.477 & 3.38 \\
  2.0 & Midpoint & 10 & 0.340 & 0.525 & 0.800 & 0.686 & 0.475 & 3.13 \\
  1.0 & Midpoint & 10 & 0.288 & 0.461 & 0.807 & 0.929 & 0.357 & 2.64 \\
  3.0 & Euler & 50 & 0.348 & 0.545 & 0.795 & 0.622 & 0.489 & 3.33 \\
  1.0 & Euler & 50 & 0.293 & 0.466 & 0.807 & 0.919 & 0.365 & 2.65 \\
  --- & Gaussian & --- & 0.011 & 0.048 & 0.466 & 2.277 & 0.009 & 0.03 \\
\midrule
\multicolumn{9}{l}{\textit{Fine-grained CFG sweep (Euler solver)}} \\
\midrule
  0.5 & Euler & 10 & 0.147 & 0.284 & 0.788 & 1.129 & --- & 2.33 \\
  1.0 & Euler & 10 & 0.348 & 0.515 & 0.796 & 0.757 & --- & 2.94 \\
  1.5 & Euler & 10 & 0.396 & 0.553 & 0.798 & 0.600 & --- & 3.36 \\
  2.0 & Euler & 10 & 0.407 & 0.580 & 0.806 & 0.526 & --- & 3.54 \\
  2.5 & Euler & 10 & 0.408 & 0.590 & 0.804 & 0.493 & --- & 3.70 \\
  3.0 & Euler & 10 & 0.410 & 0.592 & 0.802 & 0.484 & --- & 3.80 \\
  4.0 & Euler & 10 & 0.402 & 0.586 & 0.804 & 0.502 & --- & 4.01 \\
  5.0 & Euler & 10 & 0.380 & 0.576 & 0.798 & 0.548 & --- & 4.34 \\
  0.5 & Euler & 20 & 0.137 & 0.269 & 0.788 & 1.230 & --- & 2.07 \\
  1.0 & Euler & 20 & 0.334 & 0.504 & 0.794 & 0.842 & --- & 2.83 \\
  1.5 & Euler & 20 & 0.380 & 0.545 & 0.798 & 0.679 & --- & 3.19 \\
  2.0 & Euler & 20 & 0.391 & 0.569 & 0.798 & 0.600 & --- & 3.43 \\
  2.5 & Euler & 20 & 0.394 & 0.578 & 0.800 & 0.559 & --- & 3.66 \\
  3.0 & Euler & 20 & 0.398 & 0.582 & 0.800 & 0.540 & --- & 3.77 \\
  4.0 & Euler & 20 & 0.388 & 0.578 & 0.802 & 0.535 & --- & 3.97 \\
  5.0 & Euler & 20 & 0.379 & 0.567 & 0.792 & 0.551 & --- & 4.16 \\
  0.5 & Euler & 50 & 0.132 & 0.259 & 0.792 & 1.318 & --- & 1.95 \\
  1.0 & Euler & 50 & 0.324 & 0.495 & 0.794 & 0.933 & --- & 2.72 \\
  1.5 & Euler & 50 & 0.370 & 0.538 & 0.796 & 0.776 & --- & 3.14 \\
  2.0 & Euler & 50 & 0.384 & 0.563 & 0.796 & 0.699 & --- & 3.31 \\
  2.5 & Euler & 50 & 0.388 & 0.576 & 0.796 & 0.658 & --- & 3.48 \\
  3.0 & Euler & 50 & 0.392 & 0.581 & 0.792 & 0.637 & --- & 3.64 \\
  4.0 & Euler & 50 & 0.379 & 0.579 & 0.790 & 0.622 & --- & 3.86 \\
  5.0 & Euler & 50 & 0.378 & 0.570 & 0.790 & 0.626 & --- & 4.10 \\
\end{longtable}
}% end footnotesize for Table S3

% Table S4: 5 Biological Validation Datasets
\begin{table}[htbp]
\centering
\caption{The 5 biological validation datasets used for marker gene correlation, expression distribution fidelity, and Gene Ontology enrichment analysis (Section~\ref{sec:cross-dataset}). These datasets were selected to span diverse human cancer tissue types and immune cell populations. \textbf{Four of the five datasets are from the training split; only GSE183904 is held-out (Table~S2). This analysis therefore primarily evaluates in-distribution decoder-reconstructed fidelity, not out-of-distribution generalization.}}
\label{tab:s4_biovalidation}
\small
\begin{tabularx}{\textwidth}{rlllcr}
\toprule
\# & GEO Accession & Tissue & Description & Split & $n_{\text{cells}}$ \\
\midrule
  1 & GSE123902 & Lung & Lung adenocarcinoma TME & Train & 2,753 \\
  2 & GSE183904 & Gastric & Gastric cancer TME & Val & 3,000 \\
  3 & GSE123813 & Skin & Basal cell carcinoma & Train & 3,000 \\
  4 & GSE98638 & Liver & HCC T cell landscape & Train & 3,000 \\
  5 & GSE132509 & Blood & ALL bone marrow cells & Train & 2,984 \\
\midrule
  & \textbf{Total} & & & & \textbf{14,737} \\
\bottomrule
\end{tabularx}
\end{table}

% Table S5: Bootstrap 95% Confidence Intervals
\begin{table}[htbp]
\centering
\caption{Bootstrap 95\% confidence intervals (percentile method, $B{=}1{,}000$) for headline generation metrics. Resampling unit: cell types ($n{=}69$ types resampled with replacement). Classifiers (KNN, logistic regression) were trained once on real data; only the composition of evaluated types varies across bootstrap iterations. Generation configuration: CFG\,=\,1.5, Midpoint solver, 20 steps. Point estimates reflect the full (non-resampled) evaluation over all 69 types (100 generated cells per type, 6{,}900 total).}
\label{tab:s5_bootstrap_ci}
\begin{tabular}{lrrrr}
\toprule
Metric & Point Est. & 95\% CI Lower & 95\% CI Upper & Boot.\ SE \\
\midrule
KNN Top-1 Accuracy      & 0.509 & 0.421 & 0.586 & 0.043 \\
KNN Top-5 Accuracy      & 0.728 & 0.658 & 0.794 & 0.035 \\
Steering Accuracy       & 1.000 & 1.000 & 1.000 & 0.000 \\
Diversity Ratio (DivR)  & 0.969 & 0.963 & 0.974 & 0.003 \\
Linear Classifier Acc.  & 0.583 & 0.518 & 0.646 & 0.034 \\
Centroid Cosine Sim.    & 0.898 & 0.875 & 0.913 & 0.009 \\
\bottomrule
\end{tabular}
\end{table}

\begin{adjustwidth}{-\extralength}{0cm}
%} % If the paper is ``preprints'', please uncomment this parenthesis.
%\printendnotes[custom] % Un-comment to print a list of endnotes

\reftitle{References}

% Please provide the correct journal abbreviation (e.g. according to the “List of Title Word Abbreviations” http://www.issn.org/services/online-services/access-to-the-ltwa/).
% Citations and References in Supplementary files are permitted provided that they also appear in the reference list here. 

%=====================================
% References, variant A: external bibliography
%=====================================
% \bibliography{your_external_BibTeX_file}

%=====================================
% References, variant B: internal bibliography
%=====================================

% ACS format — ordered by first citation appearance in the text
\begin{thebibliography}{999}
\setlength{\itemsep}{0pt}\setlength{\parskip}{0pt}\setlength{\parsep}{0pt}

\bibitem{ref-scrnaseq-review}
Svensson,~V.; Vento-Tormo,~R.; Teichmann,~S.A. Exponential scaling of single-cell RNA-seq in the past decade. {\em Nat.\ Protoc.} {\bf 2018}, {\em 13}, 599--604.

\bibitem{ref-10x}
Zheng,~G.X.Y.; Terry,~J.M.; Belgrader,~P.; Ryvkin,~P.; Bent,~Z.W.; Wilson,~R.; Ziraldo,~S.B.; Wheeler,~T.D.; McDermott,~G.P.; Zhu,~J.; et~al. Massively parallel digital transcriptional profiling of single cells. {\em Nat.\ Commun.} {\bf 2017}, {\em 8}, 14049.

\bibitem{ref-scvi}
Lopez,~R.; Regier,~J.; Cole,~M.B.; Jordan,~M.I.; Yosef,~N. Deep generative modeling for single-cell transcriptomics. {\em Nat.\ Methods} {\bf 2018}, {\em 15}, 1053--1058.

\bibitem{ref-scvi-tools}
Gayoso,~A.; Lopez,~R.; Xing,~G.; Boyeau,~P.; Valiber Wu,~K.; Jayasuriya,~M.; Mehlman,~E.; Lancel,~M.; Regier,~J.; et~al. A Python library for probabilistic analysis of single-cell omics data. {\em Nat.\ Biotechnol.} {\bf 2022}, {\em 40}, 163--166.

\bibitem{ref-scgen}
Lotfollahi,~M.; Wolf,~F.A.; Theis,~F.J. scGen predicts single-cell perturbation responses. {\em Nat.\ Methods} {\bf 2019}, {\em 16}, 715--721.

\bibitem{ref-scdiff}
Ding,~J.; Regev,~A. Deep generative model embedding of single-cell RNA-Seq profiles on hyperspheres and hyperbolic spaces. {\em Nat.\ Commun.} {\bf 2021}, {\em 12}, 2554.

\bibitem{ref-geneformer}
Theodoris,~C.V.; Xiao,~L.; Chopra,~A.; Chaffin,~M.D.; Al Sayed,~Z.R.; Hill,~M.C.; Mantineo,~H.; Brydon,~E.M.; Zeng,~Z.; Liu,~X.S.; et~al. Transfer learning enables predictions in network biology. {\em Nature} {\bf 2023}, {\em 618}, 616--624.

\bibitem{ref-scbert}
Yang,~F.; Wang,~W.; Wang,~F.; Fang,~Y.; Tang,~D.; Huang,~J.; Lu,~H.; Yao,~J. scBERT as a large-scale pretrained deep language model for cell type annotation of single-cell RNA-seq data. {\em Nat.\ Mach.\ Intell.} {\bf 2022}, {\em 4}, 852--861.

\bibitem{ref-scfoundation}
Hao,~M.; Gong,~J.; Zeng,~X.; Liu,~C.; Guo,~Y.; Cheng,~X.; Wang,~T.; Ma,~J.; Zhang,~X.; Song,~L. Large-scale foundation model on single-cell transcriptomics. {\em Nat.\ Methods} {\bf 2024}, {\em 21}, 1481--1491.

\bibitem{ref-cellplm}
Wen,~H.; Tang,~W.; Dai,~X.; Ding,~J.; Jin,~W.; Xie,~Y.; Tang,~J. CellPLM: Pre-training of cell language model beyond single cells. In Proceedings of the 12th International Conference on Learning Representations (ICLR), Vienna, Austria, 7--11 May 2024.

\bibitem{ref-uce}
Rosen,~Y.; Brbic,~M.; Rozenblatt-Rosen,~O.; Leskovec,~J. Universal cell embeddings. In NeurIPS Workshop on New Frontiers in Learning, Control, and Dynamical Systems, New Orleans, LA, USA, 15 December 2023.

\bibitem{ref-cell2sentence}
Levine,~D.; Rizvi,~S.A.; L\'{e}vy,~S.; Pallikuth,~C.; Zhang,~X.; Chen,~E.; Zhang,~J.; Ghadermarzi,~S.; Wu,~H.; Zheng,~D.; Sontag,~D.; Pe'er,~D. Cell2Sentence: Teaching large language models the language of biology. {\em bioRxiv} {\bf 2024}, 2023.09.11.557287.

\bibitem{ref-genept}
Chen,~Y.; Bhatt,~H.; Bhatt,~A.; Liu,~C.; Langdon,~A.; Li,~C.L.; Bhatt,~P. GenePT: A simple but effective foundation model for genes and cells built from ChatGPT. {\em bioRxiv} {\bf 2024}, 2023.10.16.562533.

\bibitem{ref-clip}
Radford,~A.; Kim,~J.W.; Hallacy,~C.; Ramesh,~A.; Goh,~G.; Agarwal,~S.; Sastry,~G.; Askell,~A.; Mishkin,~P.; Clark,~J.; et~al. Learning transferable visual models from natural language supervision. In Proceedings of the 38th International Conference on Machine Learning (ICML), Virtual, 18--24 July 2021; pp.\ 8748--8763.

\bibitem{ref-siglip}
Zhai,~X.; Mustafa,~B.; Kolesnikov,~A.; Beyer,~L. Sigmoid loss for language image pre-training. In Proceedings of the IEEE/CVF International Conference on Computer Vision (ICCV), Paris, France, 2--6 October 2023; pp.\ 11975--11986.

\bibitem{ref-dit}
Peebles,~W.; Xie,~S. Scalable diffusion models with transformers. In Proceedings of the IEEE/CVF International Conference on Computer Vision (ICCV), Paris, France, 2--6 October 2023; pp.\ 4195--4205.

\bibitem{ref-flow-matching}
Lipman,~Y.; Chen,~R.T.Q.; Ben-Hamu,~H.; Nickel,~M. Flow matching for generative modeling. In Proceedings of the 11th International Conference on Learning Representations (ICLR), Kigali, Rwanda, 1--5 May 2023.

\bibitem{ref-albergo-flow}
Albergo,~M.S.; Vanden-Eijnden,~E. Building normalizing flows with stochastic interpolants. In Proceedings of the 11th International Conference on Learning Representations (ICLR), Kigali, Rwanda, 1--5 May 2023.

\bibitem{ref-rectified-flow}
Liu,~X.; Gong,~C.; Liu,~Q. Flow straight and fast: Learning to generate and transfer data with rectified flow. In Proceedings of the 11th International Conference on Learning Representations (ICLR), Kigali, Rwanda, 1--5 May 2023.

\bibitem{ref-seurat}
Stuart,~T.; Butler,~A.; Hoffman,~P.; Hafemeister,~C.; Papalexi,~E.; Mauck,~W.M.; Hao,~Y.; Stoeckius,~M.; Smibert,~P.; Satija,~R. Comprehensive integration of single-cell data. {\em Cell} {\bf 2019}, {\em 177}, 1888--1902.

\bibitem{ref-hvg}
Brennecke,~P.; Anders,~S.; Kim,~J.K.; Kolodziejczyk,~A.A.; Zhang,~X.; Prosdocimi,~V.; Dinar,~B.; Teichmann,~S.A.; Marioni,~J.C.; Heisler,~M.G. Accounting for technical noise in single-cell RNA-seq experiments. {\em Nat.\ Methods} {\bf 2013}, {\em 10}, 1093--1095.

\bibitem{ref-wilcoxon}
Wilcoxon,~F. Individual comparisons by ranking methods. {\em Biom.\ Bull.} {\bf 1945}, {\em 1}, 80--83.

\bibitem{ref-panglao}
Franz\'{e}n,~O.; Gan,~L.M.; Bj\"{o}rkegren,~J.L.M. PanglaoDB: A web server for exploration of mouse and human single-cell RNA sequencing data. {\em Database} {\bf 2019}, {\em 2019}, baz046.

\bibitem{ref-cellmarker}
Hu,~C.; Li,~T.; Xu,~Y.; Zhang,~X.; Li,~F.; Bai,~J.; Chen,~J.; Jiang,~W.; Yang,~K.; Ou,~Q.; et~al. CellMarker~2.0: An updated database of manually curated cell markers in human/mouse and web tools based on scRNA-seq data. {\em Nucleic Acids Res.} {\bf 2023}, {\em 51}, D870--D876.

\bibitem{ref-leiden}
Traag,~V.A.; Waltman,~L.; van~Eck,~N.J. From Louvain to Leiden: Guaranteeing well-connected communities. {\em Sci. Rep.} {\bf 2019}, {\em 9}, 5233.

\bibitem{ref-tabula-muris}
Tabula Muris Consortium. Single-cell transcriptomics of 20 mouse organs creates a Tabula Muris. {\em Nature} {\bf 2018}, {\em 562}, 367--372.

\bibitem{ref-batchnorm}
Ioffe,~S.; Szegedy,~C. Batch normalization: Accelerating deep network training by reducing internal covariate shift. In Proceedings of the 32nd International Conference on Machine Learning (ICML), Lille, France, 6--11 July 2015; pp.\ 448--456.

\bibitem{ref-gelu}
Hendrycks,~D.; Gimpel,~K. Gaussian error linear units (GELUs). {\em arXiv} {\bf 2016}, 1606.08415.

\bibitem{ref-dropout}
Srivastava,~N.; Hinton,~G.; Krizhevsky,~A.; Sutskever,~I.; Salakhutdinov,~R. Dropout: A simple way to prevent neural networks from overfitting. {\em J.\ Mach.\ Learn.\ Res.} {\bf 2014}, {\em 15}, 1929--1958.

\bibitem{ref-adamw}
Loshchilov,~I.; Hutter,~F. Decoupled weight decay regularization. In Proceedings of the 7th International Conference on Learning Representations (ICLR), New Orleans, LA, USA, 6--9 May 2019.

\bibitem{ref-cosine-schedule}
Loshchilov,~I.; Hutter,~F. SGDR: Stochastic gradient descent with warm restarts. In Proceedings of the 5th International Conference on Learning Representations (ICLR), Toulon, France, 24--26 April 2017.

\bibitem{ref-ot-cfm}
Tong,~A.; Malkin,~N.; Huguet,~G.; Zhang,~Y.; Rector-Brooks,~J.; Fatras,~K.; Wolf,~G.; Bengio,~Y. Improving and generalizing flow-based generative models with minibatch optimal transport. {\em Trans.\ Mach.\ Learn.\ Res.} {\bf 2024}.

\bibitem{ref-sd3}
Esser,~P.; Kulal,~S.; Blattmann,~A.; Entezari,~R.; M\"{u}ller,~J.; Saini,~H.; Levi,~Y.; Lorenz,~D.; Sauer,~A.; Boesel,~F.; et~al. Scaling rectified flow transformers for high-resolution image synthesis. In Proceedings of the 41st International Conference on Machine Learning (ICML), Vienna, Austria, 21--27 July 2024.

\bibitem{ref-cfg}
Ho,~J.; Salimans,~T. Classifier-free diffusion guidance. In Proceedings of the NeurIPS 2021 Workshop on Deep Generative Models and Downstream Applications, Virtual, 13 December 2021.

\bibitem{ref-ema}
Polyak,~B.T.; Juditsky,~A.B. Acceleration of stochastic approximation by averaging. {\em SIAM J.\ Control Optim.} {\bf 1992}, {\em 30}, 838--855.

\bibitem{ref-ari}
Hubert,~L.; Arabie,~P. Comparing partitions. {\em J.\ Classif.} {\bf 1985}, {\em 2}, 193--218.

\bibitem{ref-nmi}
Strehl,~A.; Ghosh,~J. Cluster ensembles---A knowledge reuse framework for combining multiple partitions. {\em J.\ Mach.\ Learn.\ Res.} {\bf 2002}, {\em 3}, 583--617.

\bibitem{ref-fid}
Heusel,~M.; Ramsauer,~H.; Unterthiner,~T.; Nessler,~B.; Hochreiter,~S. GANs trained by a two time-scale update rule converge to a local Nash equilibrium. {\em Adv.\ Neural Inf.\ Process.\ Syst.} {\bf 2017}, {\em 30}, 6626--6637.

\bibitem{ref-umap}
McInnes,~L.; Healy,~J.; Melville,~J. UMAP: Uniform manifold approximation and projection for dimension reduction. {\em arXiv} {\bf 2018}, 1802.03426.

\bibitem{ref-scanpy}
Wolf,~F.A.; Angerer,~P.; Theis,~F.J. SCANPY: Large-scale single-cell gene expression data analysis. {\em Genome Biol.} {\bf 2018}, {\em 19}, 15.

\bibitem{ref-integration-benchmark}
Luecken,~M.D.; Buttner,~M.; Chaichoompu,~K.; Danese,~A.; Interlandi,~M.; Mueller,~M.F.; Strobl,~D.C.; Zappia,~L.; Dugas,~M.; Colome-Tatche,~M.; Theis,~F.J. Benchmarking atlas-level data integration in single-cell genomics. {\em Nat.\ Methods} {\bf 2022}, {\em 19}, 41--50.

\bibitem{ref-clop-dit}
Fu,~Z. CLOP-DiT: Structured-Metadata-Conditioned Single-Cell Latent Generation via Contrastive Language-Omics Pretraining and Diffusion Transformers. {\em GitHub repository} {\bf 2025}. Available online: https://github.com/zeyufu/CLOP-DiT (accessed on 7 March 2026).

\end{thebibliography}

% If authors have biography, please use the format below
%\section*{Short Biography of Authors}
%\bio
%{\raisebox{-0.35cm}{\includegraphics[width=3.5cm,height=5.3cm,clip,keepaspectratio]{Definitions/author1.pdf}}}
%{\textbf{Firstname Lastname} Biography of first author}
%
%\bio
%{\raisebox{-0.35cm}{\includegraphics[width=3.5cm,height=5.3cm,clip,keepaspectratio]{Definitions/author2.pdf}}}
%{\textbf{Firstname Lastname} Biography of second author}

% For the MDPI journals use author-date citation, please follow the formatting guidelines on http://www.mdpi.com/authors/references
% To cite two works by the same author: \citeauthor{ref-journal-1a} (\citeyear{ref-journal-1a}, \citeyear{ref-journal-1b}). This produces: Whittaker (1967, 1975)
% To cite two works by the same author with specific pages: \citeauthor{ref-journal-3a} (\citeyear{ref-journal-3a}, p. 328; \citeyear{ref-journal-3b}, p.475). This produces: Wong (1999, p. 328; 2000, p. 475)

%%%%%%%%%%%%%%%%%%%%%%%%%%%%%%%%%%%%%%%%%%
%% for journal Sci
%\reviewreports{\\
%Reviewer 1 comments and authors’ response\\
%Reviewer 2 comments and authors’ response\\
%Reviewer 3 comments and authors’ response
%}
%%%%%%%%%%%%%%%%%%%%%%%%%%%%%%%%%%%%%%%%%%
\PublishersNote{}
%\isPreprints{}{% This command is only used for ``preprints''.
\end{adjustwidth}
%} % If the paper is ``preprints'', please uncomment this parenthesis.
\end{document}

